% NOTE (buck_obs_apt, 29 Oct 2026): this memo is Buck_disturb_obs's and describes the base
% controller only (the per-cycle loop, unchanged here with TOC and learning off). The TOC pulse
% pair and the learning paths added in stage 2 are described in README.md (part 1).
\documentclass[11pt]{article}

\usepackage[margin=1in]{geometry}
\usepackage[T1]{fontenc}
\usepackage{lmodern}
\usepackage{microtype}
\usepackage{amsmath,amssymb}
\usepackage{booktabs}
\usepackage{array}
\usepackage{tabularx}
\usepackage{xltabular}
\usepackage{enumitem}
\usepackage{float}
\usepackage[hidelinks]{hyperref}

% Code identifiers with underscores, typeset verbatim.
\newcommand{\code}[1]{\texttt{\detokenize{#1}}}
% Marks the option that was chosen.
\newcommand{\chosen}{\hfill$\checkmark$\,\textbf{selected}}
\newcolumntype{L}{>{\raggedright\arraybackslash}X}

\setlist[description]{font=\normalfont\bfseries, leftmargin=1.5em}
\setlist[enumerate,1]{label=(\alph*), nosep, leftmargin=2em}
\setlist[enumerate,2]{label=(\alph*), nosep, leftmargin=2em}

\setlength{\emergencystretch}{2em}

\begin{document}

% ============================================================
\begin{center}
{\Large\bfseries Memorandum}
\end{center}

\noindent
\begin{tabularx}{\linewidth}{@{}l L@{}}
\textbf{To:}   & Project owner, \code{Buck_Current_Mode} \\
\textbf{From:} & Claude Code (firmware implementation session) \\
\textbf{Date:} & 26 September 2026 (updated 28 October 2026) \\
\textbf{Re:}   & Questions and answers on the NUCLEO-G474RE
                 average-current-mode buck firmware \\
\end{tabularx}

\medskip\hrule\medskip

% ============================================================
\section{Purpose}
% ============================================================

The specification ended with the instruction to ask about anything
unclear instead of deciding. While implementing the firmware I asked
sixteen questions in four rounds. This memo records each question,
the options offered, your answer, and how the answer appears in the
code. It also lists the implementation choices I made without
asking, and the questions raised after the implementation that
still need your decision.

The firmware builds without warnings in the Makefile and in both
CubeIDE configurations, and the host tests and the closed-loop
simulation pass. Since 27 September it also runs on the board with the
power stage (Section~\ref{sec:followup}); the two-sample
$v_{\mathrm{out}}$ acquisition added later that day ran on the board in
the Buck\_disturb\_obs build you ran on 27 October (below). The copy
BUCK\_Adpt\_Learning (24 October, Section~\ref{sec:followup1024})
has not run on the board. This copy, Buck\_disturb\_obs (25 October,
Section~\ref{sec:followup1025obs}), tripped fault~5 on the first pass
after arming on its first board run that evening; the fix
(Section~\ref{sec:followup1025f}, bin e5c7798f\ldots) ran on the board on
27 October: you reported that it works (no data from that run were read).
The clock optimisation of 28 October (Section~\ref{sec:followup1028}; bin
5b99a0d5\ldots, the applied 2f144a53\ldots{} with the observer's prediction
pinned in the post pass) gives the same results bit for bit but for one
exception, 5 debugger-stall scenarios in which the 25 October build trips
its own fault~5 (row~9 there). It has not run on the board yet.

% ============================================================
\section{Summary of Answers}
% ============================================================

\begin{table}[H]
\centering
\caption{All sixteen questions. Answers in quotation marks were typed
by you rather than picked from the options.}
\label{tab:summary}
\footnotesize
\renewcommand{\arraystretch}{1.15}
\begin{tabularx}{\linewidth}{@{}r>{\raggedright\arraybackslash}p{0.25\linewidth}LL@{}}
\toprule
\textbf{\#} & \textbf{Question} & \textbf{Answer} & \textbf{Where in the code} \\
\midrule
1  & Loop structure; which $u[k]$ enters $S[k]$
   & Cascaded; $u[k]=i_{\mathrm{ref}}[k]$, the voltage-PI output in A
     (24 October: $S$ uses the duty change $\Delta d$ instead; 25 October:
     $S=a e_v^2+b e_v+c$, no second variable)
   & \code{buck_ctrl_fast()}; \code{out.i_ref} \\
2  & Feedback of the inner current loop
   & Corrected estimate $\hat{\bar{i}}_L[k]$
   & current-PI error $i_{\mathrm{ref}}-\hat{\bar{i}}_L[k]$ \\
3  & Acquisition of $\bar{i}_{L,\mathrm{m}}[k]$
   & 4 samples per period, averaged in hardware (ADC oversampling)
   & HRTIM Master CMP1--4 $\to$ ADC2, 4$\times$ oversampling \\
4  & Firmware style and clock
   & Register-level CMSIS, 168\,MHz from HSI
   & \code{clock_init()}; HRTIM period 8960 \\
\midrule
5  & Source of $V_{\mathrm{in}}$
   & Constant 6\,V
   & \code{BUCK_VIN_V} \\
6  & A0 scaling and $V_{\mathrm{ref}}$
   & A0 direct, $V_{\mathrm{ref}}=1000$ counts
   & \code{tune.vref_counts} \\
7  & A1 output at 0\,A
   & 1.65\,V (bidirectional); corrected later to 0\,V at 0\,A
   & \code{tune.isense_zero} $=0$ (was 2048) \\
8  & Inductor resistance $r_L$
   & 114\,m$\Omega$
   & \code{BUCK_RL_OHM} \\
\midrule
9  & PI gains and limits
   & ``use textbook 45 phase margin''
   & \code{tools/pi_design.py}; crossovers in Q13 \\
10 & Correction factor $\alpha_I$
   & $\tau_c=50\,\mu\mathrm{s}$, giving $\alpha_I=0.03278$
   & \code{BUCK_ALPHA_I}; \code{tune.alpha_i} \\
11 & Trigger $\rho$ and $S_{\mathrm{th}}$
   & $\rho=0$, $S_{\mathrm{th}}=0.01$; $S_{\mathrm{th}}$ changed later to 1;
     $\Delta d$ for $\Delta u$ on 24 October; $\rho$ removed on 25 October
   & \code{tune.s_th}; \code{tune.score_a}, \code{_b}, \code{_c}
     (\code{tune.rho} until 25 October) \\
12 & Scaling of $\hat{I}$ on A2
   & 0--3.3\,V $\leftrightarrow$ 0--10\,A
   & \code{BUCK_DAC_FULL_SCALE_A} \\
\midrule
13 & Crossover pair for the $45^\circ$ design
   & $f_{ci}=60\,\mathrm{kHz}$, $f_{cv}=6\,\mathrm{kHz}$
   & \code{BUCK_KP_I}, \code{BUCK_KI_I}, \code{BUCK_KP_V}, \code{BUCK_KI_V} \\
14 & Dead time in $T_{\mathrm{on}}[k]$
   & ``consider deadtime as off time''
   & CMP1 $=D\cdot 8960+272$ \\
15 & Volts per A0 count
   & $3.3\,\mathrm{V}/4096=805.7\,\mu\mathrm{V}$
   & \code{BUCK_V_PER_COUNT} \\
16 & Start-up of the gates
   & Off until a debugger command
   & \code{g_buck.command} \\
\bottomrule
\end{tabularx}
\end{table}

% ============================================================
\section{Round 1: Control Architecture}
% ============================================================

\subsection*{Q1.\ Loop structure and the meaning of $u[k]$}
\begin{description}
\item[Why asked.] The specification calls for average-current-mode
control but describes a single PI with output $u[k]$, and
$\Delta u[k]$ enters the trigger score $S[k]$.
\item[Options offered.]\leavevmode
  \begin{enumerate}
  \item Cascaded, $u={}$voltage-PI output: the outer voltage PI gives
        the average-current reference $i_{\mathrm{ref}}$; the inner
        current PI gives the duty; $\Delta u$ is in amperes.\chosen
  \item Cascaded, $u={}$current-PI output: $\Delta u$ is in duty units.
  \item Single PI from the voltage error to duty (voltage mode; the
        estimator only feeds the DAC).
  \end{enumerate}
\item[Implemented as.] $u[k]=i_{\mathrm{ref}}[k]$ in amperes, and
$\Delta u[k]=i_{\mathrm{ref}}[k]-i_{\mathrm{ref}}[k-1]$ in $S[k]$.
\item[Change.] On 24 October (BUCK\_Adpt\_Learning) $S[k]$ uses the duty
change $\Delta d[k]=d[k]-d[k-1]$ of the current PI instead of $\Delta u$
(Section~\ref{sec:followup1024}). On 25 October the $\Delta d$ term was
removed: $S=a e_v^2+b e_v+c$ (same section).
\end{description}

\subsection*{Q2.\ Feedback of the inner current loop}
\begin{description}
\item[Options offered.]\leavevmode
  \begin{enumerate}
  \item The measured cycle average $\bar{i}_{L,\mathrm{m}}[k]$ from A1;
        the estimator is used only for $\hat{I}$ and the trigger.
  \item The measurement-corrected estimate $\hat{\bar{i}}_L[k]$ of the
        specification's corrected-average equation.\chosen
  \end{enumerate}
\item[Implemented as.] The current PI acts on
$i_{\mathrm{ref}}[k]-\hat{\bar{i}}_L[k]$.
\end{description}

\subsection*{Q3.\ Acquisition of the measured cycle average}
\begin{description}
\item[Why asked.] The specification defines
$\bar{i}_{L,\mathrm{m}}[k]$ but not how it is sampled from the
650\,kHz sensor on A1.
\item[Options offered.]\leavevmode
  \begin{enumerate}
  \item Four samples per period at $T/8$, $3T/8$, $5T/8$, $7T/8$,
        averaged by ADC oversampling; no CPU cost; window aligned to
        the switching cycle.\chosen
  \item One sample at mid on-time (sensitive to the sensor's
        $\approx 245\,\mathrm{ns}$ lag).
  \item One sample plus a computed ripple correction, as in
        \code{Buck_Voltage_Mode}.
  \end{enumerate}
\item[Implemented as.] HRTIM Master CMP1--CMP4 drive ADC trigger~3.
ADC2 runs triggered $4\times$ oversampling with no shift, so one
result per cycle holds the sum of the four samples.
\end{description}

\subsection*{Q4.\ Firmware style and clock}
\begin{description}
\item[Options offered.]\leavevmode
  \begin{enumerate}
  \item Register-level CMSIS, 168\,MHz from HSI: HRTIM $\times 32$
        gives exactly $8960$ ticks $=600.000\,\mathrm{kHz}$; interrupt
        handler placed in CCM RAM.\chosen
  \item Register-level CMSIS, 168\,MHz from the 24\,MHz HSE (crystal
        accuracy; the board's HSE routing must be checked).
  \item HAL/LL generated by CubeMX (heavier code path; 170\,MHz cannot
        give an exact 600\,kHz period).
  \end{enumerate}
\item[Implemented as.] HSI16 $/4\times 84$ (336\,MHz VCO), R$/2$ gives
168\,MHz for the CPU and HRTIM, P$/6$ gives 56\,MHz for the ADCs. The
period is 280 CPU clocks. The control code runs from CCM RAM as offered,
but as a polling loop instead of an interrupt handler. This departs from
the option text and was not asked; see C1.
\end{description}

% ============================================================
\section{Round 2: Sensing and Model Parameters}
% ============================================================

\subsection*{Q5.\ Source of $V_{\mathrm{in}}[k]$}
\begin{description}
\item[Why asked.] The estimator needs $V_{\mathrm{in}}[k]$, but no
input-voltage pin is assigned.
\item[Options offered.]\leavevmode
  \begin{enumerate}
  \item Constant 6\,V.\chosen
  \item Add a $V_{\mathrm{in}}$ sense pin.
  \end{enumerate}
\item[Implemented as.] $V_{\mathrm{in}}=6\,\mathrm{V}$ at compile time
(\code{BUCK_VIN_V}).
\end{description}

\subsection*{Q6.\ A0 scaling and the reference $V_{\mathrm{ref}}$}
\begin{description}
\item[Options offered.]\leavevmode
  \begin{enumerate}
  \item A0 $=v_{\mathrm{out}}$ directly, $V_{\mathrm{ref}}=1.0\,\mathrm{V}$.
  \item A0 directly, $V_{\mathrm{ref}}=1000$ counts
        ($\approx 0.8\,\mathrm{V}$), as in
        \code{Buck_Voltage_Mode}.\chosen
  \item Divider of $0.4128\,\mathrm{V/V}$, $V_{\mathrm{ref}}=1.0\,\mathrm{V}$.
  \end{enumerate}
\item[Implemented as.] \code{tune.vref_counts} $=1000$, which is
$0.8057\,\mathrm{V}$ with the scale of Q15. It can be changed live.
\end{description}

\subsection*{Q7.\ A1 output at zero inductor current}
\begin{description}
\item[Options offered.]\leavevmode
  \begin{enumerate}
  \item 0\,V at 0\,A (unidirectional; negative current clips).
  \item 1.65\,V at 0\,A (bidirectional).\chosen
  \item Calibrate automatically at start-up with the gates off.
  \end{enumerate}
\item[Implemented as.] Zero at 2048 counts and
$20.14\,\mathrm{mA}$ per count. The zero is live-tunable as
\code{tune.isense_zero} and the gain as \code{tune.isense_gain}, so the
zero can be trimmed by hand from the debugger. There is no automatic
calibration.
\item[Correction.] On 27 September you corrected this answer: A1 reads
0\,V at 0\,A, as in option~(a). The default zero is now 0 counts
(\code{BUCK_ISENSE_ZERO_V} $=0$), and negative inductor current reads as
0\,A. In simulation the clipping changes nothing from about 1\,A up and
makes the measured average read high at lighter load (about
$+0.1\,\mathrm{A}$ at 0.5\,A).
\end{description}

\subsection*{Q8.\ Inductor series resistance $r_L$}
\begin{description}
\item[Options offered.]\leavevmode
  \begin{enumerate}
  \item $114\,\mathrm{m}\Omega$, the DC fit from
        \code{Buck_Voltage_Mode} (includes FET and path
        resistance).\chosen
  \item $0\,\Omega$ ($\theta=0$).
  \end{enumerate}
\item[Implemented as.] $\theta=r_L/L=3.45\times 10^{5}\,\mathrm{s}^{-1}$,
so $\theta T_{\mathrm{sw}}=0.58<2$. Every Pad\'e $a(\tau)$ is therefore
positive and $\kappa_I$ is always defined.
\end{description}

% ============================================================
\section{Round 3: Gains, Correction, Trigger and DAC}
% ============================================================

\subsection*{Q9.\ PI gains and limits}
\begin{description}
\item[Options offered.]\leavevmode
  \begin{enumerate}
  \item Reuse the gains of \code{Buck_Voltage_Mode}: voltage PI
        $K_p=30.8\,\mathrm{A/V}$, $K_i=1.155\times10^{6}\,\mathrm{A/(V\,s)}$;
        current PI $K_p=0.0279\,\mathrm{A^{-1}}$,
        $K_i=1570\,\mathrm{(A\,s)^{-1}}$; limits $i_{\mathrm{ref}}$
        $-2\ldots10\,\mathrm{A}$ and duty 10--50\,\%.
  \item Enter values.
  \end{enumerate}
\item[Answer.] Typed: ``use textbook 45 phase margin''.
\item[Implemented as.] Both PIs are designed for $45^\circ$ phase
margin by \code{tools/pi_design.py}. The crossover frequencies were
asked in Q13. The answer did not cover the limits, so they were carried
over from \code{Buck_Voltage_Mode}; they are still open (O2 in
Section~\ref{sec:open}).
\end{description}

\subsection*{Q10.\ Measurement-correction factor $\alpha_I$}
\begin{description}
\item[Options offered.]\leavevmode
  \begin{enumerate}
  \item $\tau_c=50\,\mu\mathrm{s}$, so
        $\alpha_I=1-e^{-T_{\mathrm{sw}}/\tau_c}\approx 0.0328$.\chosen
  \item Fixed $\alpha_I=0.05$.
  \item Fixed $\alpha_I=0.1$.
  \end{enumerate}
\item[Implemented as.] $\alpha_I=0.03278$, a constant because
$T_{\mathrm{sw}}$ is fixed. It is evaluated at compile time and can be
changed live (\code{tune.alpha_i}).
\end{description}

\subsection*{Q11.\ Trigger parameters}
\begin{description}
\item[Options offered.]\leavevmode
  \begin{enumerate}
  \item $\rho=0$, $S_{\mathrm{th}}=0.01$: fires at about
        $|e_v|\ge 100\,\mathrm{mV}$ or $|\Delta u|\ge 0.1\,\mathrm{A}$.\chosen
  \item Enter values.
  \end{enumerate}
\item[Implemented as.] Both values are live-tunable (\code{tune.rho},
\code{tune.s_th}); $e_v$ is in volts and $\Delta u$ in amperes.
\item[Change.] On 27 September you set the default $S_{\mathrm{th}}$ to 1
(Section~\ref{sec:followup}). On 24 October the second term became the
duty change $\Delta d$ (a plain number; \code{du_norm} kept its name). On
25 October the score became $S=a e_v^2+b e_v+c$ (\code{tune.score_a},
\code{_b}, \code{_c}; defaults 1, 0, 0): $\rho$, the denominators and
$\Delta d$ are gone, and $S_{\mathrm{th}}$ sets D4 only
(Section~\ref{sec:followup1024}).
\end{description}

\subsection*{Q12.\ Scaling of $\hat{I}$ on the A2 DAC}
\begin{description}
\item[Options offered.]\leavevmode
  \begin{enumerate}
  \item $40\,\mathrm{mV/A}$ with 1.65\,V at 0\,A, the same scale as A1.
  \item $0$--$3.3\,\mathrm{V}\leftrightarrow 0$--$10\,\mathrm{A}$
        ($0.33\,\mathrm{V/A}$, unipolar).\chosen
  \item 1.65\,V at 0\,A, $\pm 10\,\mathrm{A}$ ($0.165\,\mathrm{V/A}$).
  \end{enumerate}
\item[Implemented as.] Full scale 10\,A. Negative estimates clip at 0\,V
and estimates above 10\,A clip at 3.3\,V.
\end{description}

% ============================================================
\section{Round 4: Follow-up Questions}
% ============================================================

\subsection*{Q13.\ Crossover frequencies for the $45^\circ$ design}
\begin{description}
\item[Why asked.] Q9 fixes the phase margin but not the crossovers.
The design is discrete-time and includes the 2-cycle
sample-to-duty delay and $r_L$. It is evaluated at no load, the worst
case for the voltage loop.
\item[Options offered.] See Table~\ref{tab:crossover}.
\item[Implemented as.] Option~(c), with the exact values
$K_{p,i}=0.0237243\,\mathrm{A^{-1}}$,
$K_{i,i}=3338.07\,\mathrm{(A\,s)^{-1}}$,
$K_{p,v}=9.3448\,\mathrm{A/V}$ and
$K_{i,v}=285\,353\,\mathrm{A/(V\,s)}$. The current loop has a phase
margin of $45.0^\circ$ and a gain margin of 3.2\,dB. The voltage loop
has $45.0^\circ$ (no load) to $61.8^\circ$ (3\,A) of phase margin and
a gain margin of 15.6\,dB. All closed-loop poles satisfy
$|z|\le 0.976$ over 0--3\,A. The gains are live-tunable.
\end{description}

\begin{table}[H]
\centering
\caption{Crossover options offered in Q13. Current-PI gains in
$\mathrm{A^{-1}}$ and $\mathrm{(A\,s)^{-1}}$; voltage-PI gains in
$\mathrm{A/V}$ and $\mathrm{A/(V\,s)}$. Blank cells were not stated in the option text.}
\label{tab:crossover}
\small
\renewcommand{\arraystretch}{1.25}
\begin{tabular}{@{}lllllllll@{}}
\toprule
 & $f_{ci}$ & $f_{cv}$ & $K_{p,i}$ & $K_{i,i}$ & GM$_i$
 & $K_{p,v}$ & $K_{i,v}$ & GM$_v$ \\
\midrule
(a) & 30\,kHz & 6\,kHz & 0.00089 & 3809 & 7.7\,dB
    & 9.10  & $2.62\times10^{5}$ & 13.1\,dB \\
(b) & 40\,kHz & 8\,kHz & 0.00856 & 4758 & 6.4\,dB
    & 11.68 & $4.55\times10^{5}$ & 13.4\,dB \\
(c)$\;\checkmark$ & 60\,kHz & 6\,kHz & 0.0237 & 3338 & 3.2\,dB
    & 9.34  & $2.85\times10^{5}$ & \\
(d) & 30\,kHz & 3\,kHz & & & & 4.39 & $7.27\times10^{4}$ & 19.6\,dB \\
\bottomrule
\end{tabular}
\end{table}

\subsection*{Q14.\ Dead time in the estimator's on-time}
\begin{description}
\item[Why asked.] The estimator needs $T_{\mathrm{on}}[k]$, and the
specification does not say how the $50\,\mathrm{ns}$ dead times are
counted.
\item[Options offered.]\leavevmode
  \begin{enumerate}
  \item Commanded on-time minus 50\,ns, with
        $T_{\mathrm{off}}=T_{\mathrm{sw}}-T_{\mathrm{on}}$.
  \item Commanded on-time as is, ignoring the dead time.
  \end{enumerate}
\item[Answer.] Typed: ``consider deadtime as off time''.
\item[Implemented as.] The duty $D$ is the physical high-side on-time
fraction. Both dead times are off-time, in which the inductor sees
$-v_{\mathrm{out}}$. The compare is CMP1 $=D\cdot 8960+272$ ticks,
$T_{\mathrm{on}}=D\,T_{\mathrm{sw}}$ and
$T_{\mathrm{off}}=T_{\mathrm{sw}}-T_{\mathrm{on}}$. This assumption is
the source of the load-estimate bias in O1.
\end{description}

\subsection*{Q15.\ Volts per A0 count}
\begin{description}
\item[Options offered.]\leavevmode
  \begin{enumerate}
  \item Nominal $3.3\,\mathrm{V}/4096=805.7\,\mu\mathrm{V}$, so
        $V_{\mathrm{ref}}=0.8057\,\mathrm{V}$.\chosen
  \item $792.6\,\mu\mathrm{V}$, from the bench measurement in
        \code{Buck_Voltage_Mode}.
  \end{enumerate}
\end{description}

\subsection*{Q16.\ Start-up of the gates}
\begin{description}
\item[Options offered.]\leavevmode
  \begin{enumerate}
  \item Off until a debugger command: PWM and ADC timing run, the gates
        stay inactive until \code{command = 1} is written, then a 10\,ms
        reference soft-start; writing 0 stops.\chosen
  \item Automatic start after ADC warm-up, with a 10\,ms soft-start.
  \end{enumerate}
\item[Implemented as.] The variable is \code{g_buck.command} (the
option text said \code{g_ctrl.command}). The reference starts at the
lower of the measured $v_{\mathrm{out}}$ and $V_{\mathrm{ref}}$, then
moves at a fixed slew of $0.8057\,\mathrm{V}$ per 10\,ms (about
$0.08\,\mathrm{V/ms}$). The ramp therefore takes 10\,ms from 0\,V to the
default $V_{\mathrm{ref}}$, and less from a pre-biased output.
\end{description}

% ============================================================
\section{Choices Made Without Asking}
% ============================================================

These choices were made during implementation without a question.
None of them changes the specified algorithm. C1 departs from the text
of the option chosen in Q4.

\begin{enumerate}[label=C\arabic*., leftmargin=2.5em]
\item \textbf{Polling loop instead of an interrupt handler.} The option
chosen in Q4 said the handler would be an interrupt in CCM RAM. A first
interrupt-driven version was estimated at about 290 of the 280 clocks
per period. The loop now runs from CCM SRAM and polls the same
end-of-conversion flag of ADC1. It is still one update per cycle, and it
saves about 50 clocks of exception overhead.
\item \textbf{Pre-computed Pad\'e coefficients.} $a_{\mathrm{on}}$,
$b_{\mathrm{on}}$, $a_{\mathrm{off}}$, $b_{\mathrm{off}}$ and
$A_I/H_I$ depend only on CMP1, which takes 3585 values. They are
computed once at start-up with the specification's formulas (70\,KB
of RAM), so the results are identical to computing them every cycle.
\item \textbf{Controller state in CPU registers} while running; it is
written back to \code{g_buck.st} on a stop.
\item \textbf{Dead time rounded up} to 68 counts $\times 0.744\,\mathrm{ns}
= 50.6\,\mathrm{ns}$, so it is never shorter than 50\,ns.
\item \textbf{Sampling instants of $v_{\mathrm{out}}$:} two samples
per cycle since 27 September (Section~\ref{sec:followup}), triggered
about 457 and 100\,ns before the period boundary, in the quiet
low-side interval. Their mean is $v_{\mathrm{out}}[k+1]$: the end
value of cycle $k$ and the start value of cycle $k+1$. (Originally
one sample, triggered 100\,ns before the boundary.)
\item \textbf{PI form and anti-windup:}
$C(z)=K_p+K_iT\,z/(z-1)$ with integrator tracking of the output limit.
\item \textbf{Faults latch} with both gates off (4: ADC overrun or
missing current result while running, or no ADC1 sample at all;
5: CMP1 deadline missed; 1, 2, 3, 6: start-up
timeouts; 8: CPU fault).
\item \textbf{ADC bulb sampling} on both ADCs: the input is tracked
from the end of the previous conversion and held 2.5 ADC clocks after
the trigger (RM0440 Figure~93).
\code{-DBUCK_ADC_BULB=0} switches back to standard sampling.
\end{enumerate}

% ============================================================
\section{Open Questions}
\label{sec:open}
% ============================================================

These were raised in the implementation report and have not yet been
answered.

\begin{enumerate}[label=O\arabic*., leftmargin=2.5em]
\item \textbf{Bias of the load estimate $\hat{I}$.} In a simulation
with real body-diode conduction during the dead times, $\hat{I}$ reads
$-1.35\,\mathrm{A}$ in error at 1\,A load and $+0.35\,\mathrm{A}$ at
3\,A. When the simulated dead time behaves exactly as the model assumes
(Q14), the error is only 5--8\,mA, so the code matches the
specification. At light load the valley current is negative, and the
high-side body diode conducts during the rising-edge dead time, which
the model counts as off-time. With $A_I\approx 0.56$ per cycle,
$\alpha_I=0.033$ is too weak to remove the offset. Regulation is not
affected, and the measured average is accurate to about
$\pm 0.04\,\mathrm{A}$. Options:
  \begin{enumerate}
  \item accept the bias as specified;
  \item raise $\alpha_I$ (0.3 halves it to $-0.73$ / $+0.18\,\mathrm{A}$;
        live-tunable, no code change);
  \item add a dead-time term that depends on the current direction
        (not in the specification).
  \end{enumerate}
\item \textbf{Limits.} The specification and the answers give none,
so these were carried over from \code{Buck_Voltage_Mode}: CMP1 between
10 and 50\,\% of the period (physical $D$ between 6.96 and 46.96\,\%),
and $i_{\mathrm{ref}}$ between $-2$ and $10\,\mathrm{A}$. Keep them, or
give values.
\item \textbf{Arming mid-period.} Writing \code{command = 1} enables
the gates part-way through a period, so the low side conducts for the
rest of that period. The estimator models that period with
$T_{\mathrm{on}}=0$. With a pre-biased output this causes a brief
negative current. Keep this, or synchronise arming to a period
boundary.
\end{enumerate}

\paragraph{Noted risks (no decision needed unless you want a change).}
\begin{itemize}[nosep]
\item The current-loop gain margin is 3.2\,dB at the chosen 60\,kHz
      crossover; option~(a) of Q13 gives 7.7\,dB.
\item Bulb sampling combined with triggered oversampling has run on
      the board on ADC2; ADC1 uses it too since the two-sample
      $v_{\mathrm{out}}$ change, which ran on the board in the build you
      ran on 27 October (reported working; no data were read; fallback
      in C8).
\item The 600\,kHz frequency carries the HSI tolerance (Q4).
\item $V_{\mathrm{in}}$ is a constant (Q5). At a real 6.5\,V, the
      simulated $\hat{I}$ error is $-1.85$ / $-0.40\,\mathrm{A}$.
\item A 1\,A $\to$ 3\,A load step dips about 200\,mV. This is expected
      from the 6\,kHz voltage crossover with 270\,$\mu$F.
\end{itemize}

% ============================================================
\section{Follow-up Requests, 27 September}
\label{sec:followup}
% ============================================================

{\small
\renewcommand{\arraystretch}{1.25}
\begin{xltabular}{\linewidth}{@{}>{\raggedright\arraybackslash}p{0.27\linewidth}L@{}}
\caption{Requests and questions after the first bench runs.}
\label{tab:followup}\\
\toprule
\textbf{Request or question} & \textbf{Answer and result} \\
\midrule
\endfirsthead
\toprule
\textbf{Request or question} & \textbf{Answer and result} \\
\midrule
\endhead
Fault 5 on the board; optimise the loop to fit the timing budget.
  & The pass was at or above the 280-clock period on average, so each
    pass started later until the CMP1 write missed its deadline. The loop
    was restructured without changing the algorithm or the per-cycle
    outputs: separate multiply and add instead of VMLA, telemetry and
    tunables every 16th pass, peripheral bases in a register, one branch
    for the exit tests, and the timer read moved next to the write.
    Static count: 229--235 $\to$ 177 core clocks per pass. Board result:
    no fault, \code{busy_max} $=220$. \\
A1 reads 0\,V at 0\,A, not 1.65\,V.
  & Corrected (Q7). Negative inductor current reads 0\,A. In simulation
    this changes nothing from about 1\,A up; at lighter load the measured
    average reads high (about $+0.1\,\mathrm{A}$ at 0.5\,A and
    $+0.2\,\mathrm{A}$ at 0.1\,A). \\
Set $S_{\mathrm{th}}$ to 1.
  & Default changed to 1. With $\rho=0$, $D4$ fires when
    $e_v^2+\Delta u^2\ge 1$, e.g.\ $|\Delta u|\ge 1\,\mathrm{A}$ in one
    cycle, or an overvoltage of 1\,V; the simulated 1\,A $\to$ 3\,A step
    peaks at $S\approx 0.05$, so $D4$ stays low for it. \\
How to put more weight on the estimator than on the measurement?
  & Lower $\alpha_I$ (\code{tune.alpha_i}, i.e.\ a longer $\tau_c$). In
    simulation this does not reduce the $\hat{I}$ noise at all (0.256\,A
    rms for every $\alpha_I$ from 0.28 to 0), and it increases the
    dead-time bias ($-1.35 \to -1.46\,\mathrm{A}$ at 1\,A). \\
Why is $\hat{I}$ so noisy with $\alpha_I=0$, a steady $v_{\mathrm{out}}$
and a small $\Delta u$?
  & $\hat{I}$ multiplies the change between two A0 samples by
    $C/T_{\mathrm{sw}}=162\,\mathrm{A/V}$: one A0 count is 0.13\,A. A
    few mV of cycle-to-cycle A0 noise, invisible on the scope, gives the
    observed jumps. $\alpha_I$ and $\Delta u$ do not enter that term. \\
Implement the $N$-cycle average for the load current.
  & $\hat{I}$ on A2 and in \code{out.i_load} is now the mean over the last
    $N$ cycles, which equals the charge balance over $N$ cycles; the
    intermediate A0 samples cancel, so A0 noise drops by $N$. You chose
    $N=16$ as the default; it is live-tunable (\code{tune.iload_n},
    1--16; $N=1$ is the specification). Simulated at $\pm 2$ counts of A0
    noise: 0.26 $\to$ 0.016\,A rms; a load step reaches 50\,\% / 90\,\%
    after 20 / 50\,$\mu$s. It costs 11 clocks per pass (expected
    \code{busy_max} about 230; fault near 255). This is not the same as
    sampling A0 several times within one cycle, which would give only
    $\sqrt{M}$ with $M\le 2$ at high duty. (Superseded on 27 September by
    the fixed two-stage average below.) \\
Keep optimising the clock cycles; ADC1 and ADC2 may both sample $v$; use
multiple agents to verify.
  & Two rounds of parallel agents searched for savings and three independent
    agents verified the combination (1190 emulated scenarios on the linked
    machine code, host tests, independent cycle counts, a silicon review).
    Result, bit-identical outputs: 166--190 $\to$ 149 core clocks per
    steady pass, 183--211 $\to$ 154 during the soft-start (the main
    finding: GCC had optimised the loop for size). Sampling $v$ on both
    ADCs was evaluated and not adopted: it costs 8--11 clocks per pass and
    only averages the ADCs' own noise. Also added: NaN-safe PI limits, an
    earlier arming window, \texttt{make check}, and a boot-time bus
    benchmark (\code{g_buck.bench}). \\
Can DMA oversample and average the ADC?
  & DMA only moves results; averaging would still cost CPU time. The
    ADC's own oversampling adds the samples in hardware, with no CPU and
    no DMA; A1 already uses it (4 samples per cycle). \\
Can ADC3, ADC4 and ADC5 sample $v_{\mathrm{out}}$ 10\,ns apart for extra
samples?
  & A0 (PA0) reaches only ADC1 and ADC2. ADC3 and ADC4 would need
    $v_{\mathrm{out}}$ jumpered to one of their pins (e.g.\ A3 = PB0,
    PB12); ADC5's pins are the gate outputs. A 10\,ns spacing cannot be set
    (the trigger is resynchronised to the 56\,MHz ADC clock, about 18\,ns),
    and samples that close average only the ADCs' own noise. Not adopted. \\
Option 1: oversample $v_{\mathrm{out}}$ twice per cycle on ADC1.
  & Two A0 samples per cycle (Timer A CMP4 and CMP3, 357\,ns apart, the
    second about 100\,ns before the period end), summed by ADC1's
    triggered $2\times$ oversampling; the loop uses their mean. 357\,ns
    rather than 268\,ns: RM0440 ignores a trigger that arrives during a
    conversion (15 ADC clocks); with the ADC and HRTIM clocks locked
    1:3, 20 clocks always leave 5 idle. Loop code unchanged apart from the scale factor
    (149 clocks per steady pass). Simulated: per-cycle $\hat{I}$ noise
    0.26 $\to$ 0.19\,A rms, 0.016 $\to$ 0.012\,A over $N=16$. The DC
    output is now within 0.4\,mV of $V_{\mathrm{ref}}$ (0.5--1.3\,mV
    above it before: the single end-of-cycle sample sat 0.7--1.1\,mV
    below the cycle average of the 3.0--3.6\,mV capacitor ripple).
    Not yet run on the board.
    RM0440 Figure~93 also showed that in bulb mode the input is held
    2.5 ADC clocks (45\,ns) after the trigger, not at it. The Live
    Expression for $v_{\mathrm{out}}$ is now
    \code{adc_v_raw*3.3/8192.0}. \\
Add live denominators for $e_v$ and $\Delta u$ in the trigger score.
  & $S=(e_v/\code{ev_norm})^2+(\Delta u/\code{du_norm})^2
    +\rho\,(e_v/\code{ev_norm})(\Delta u/\code{du_norm})$, with
    \code{tune.ev_norm} in V and \code{tune.du_norm} in A (both $>0$).
    The defaults, 1\,V and 1\,A, give the specification's score bit for
    bit. The reciprocals are formed every 64 cycles with the other
    tunables (one division for both), so no cycle divides; the extra
    multiplies and registers cost 6--8 clocks per pass (149 $\to$ 155 in
    steady state). The division on the arming path narrowed the
    late-arming margin, so \code{BUCK_ARM_LATEST} went from 5120 to 4096
    ticks (normal arming is at about 2000--2600). Four independent agents
    checked it: the machine code's score matches the formula to
    $2\times10^{-7}$ with live $\rho$ and denominators; at the defaults
    every output is bit-identical to before; 0 or NaN keeps D4 low, a tiny
    value keeps it high, and the control is never affected. New tests pin
    the sign and order of $e_v$ and $\Delta u$ and $u[k-1]$ at arming. \\
Why is $\hat{I}$ on A2 still noisy? Is $N=16$ applied? (5\,mV of
$\Delta v$ is 0.8\,A, 0.26\,V on A2.)
  & $N=16$ is applied: checked in the source, in the machine code
    (emulated DAC writes equal the 16-cycle mean on every pass, from the
    first pass after arming) and in the flashed image
    (\code{iload_n} $=16$). The arithmetic is right for $N=1$; at $N=16$
    the $\Delta v$ term falls by 16, so 5\,mV gives 0.05\,A (17\,mV on A2).
    Expected: 2.7\,mV rms (about 20\,mV p-p) per LSB of A0 noise at
    $N=16$, 44\,mV rms at $N=1$. Things $N$ does not remove: the
    estimate is negative below about 1.1\,A (dead-time bias), so A2 sits
    at 0\,V; slow content (load changes, $V_{\mathrm{in}}$ ripple at
    $-1.45$\,A/V); and analog pickup on the A2 probe. Board checks:
    compare A2 at \code{iload_n} 1 and 16, snapshot \code{out.i_load},
    measure at 2.5\,A or more (README, \emph{What to expect on A2}).
    (Superseded: \code{iload_n} was removed later that day; to compare
    with the per-cycle $\hat{I}$, rebuild with $N_1=N_2=1$.) \\
Are all $v$ in the calculations the two-sample mean?
  & Yes: the voltage PI, the estimator's prediction, $\hat{I}$ (both
    $v[k+1]$ and $v[k]$), $e_v$ and arming all use it; checked in the C
    code and in the machine code. It lowers uncorrelated ADC noise by
    $\sqrt{2}$ (0.26 $\to$ 0.19\,A per cycle); noise common to both
    samples or picked up at a fixed phase is not reduced. \\
Fault 5 on the board with the two-sample and trigger-denominator build;
fix it.
  & Leading suspect (board data not yet read): the 64th pass published
    the telemetry and re-read the tunables in one pass, with two float
    divisions and an integer division, about 300--325 core clocks (period
    280), so the next pass's CMP1 write could miss its deadline; a
    verifier found a second path, a restart whose first pass was that
    pass. Fixed: telemetry and tunables in different passes 32 apart,
    the first of them 24 passes after arming, and the load-window
    divisions removed. The split itself left D4, A2 and CMP1 unchanged
    (140\,190 emulated writes; only \code{out.i_load} could differ by up
    to 2 ulp); that comparison predates the 24-pass start and the
    two-stage average, which were checked separately (the first telemetry
    or tunables pass is the 24th after arming for every starting
    \code{loops}, all $2^{32}$ values). New
    \code{g_buck.trip_cnt} records the timer at a missed deadline; with
    \code{loops} mod 64 it identifies the pass. \\
Try a two-stage load-current average of 32 $\times$ 32 cycles; verify the
clocks with several agents.
  & Stage 1 is the previous charge-balance mean over $N_1=32$ cycles,
    stage 2 the mean of the last $N_2=32$ stage-1 values. A0 noise falls
    by $N_1\sqrt{N_2}=181$ (16 before); the $\hat{\bar{\imath}}_L$ part
    only by about 3 in simulation (7 for white noise), and it makes up most
    of the simulated 0.0023 / 0.0015\,A rms at 1 / 3\,A with $\pm2$ counts
    (0.012\,A before); a load step
    reaches 90\,\% after about 100\,$\mu$s (50 before). You chose fixed
    lengths: a live-tunable pair measured about 8 clocks more in every
    pass and brought two divisions back; fixed powers of two make the A2
    code an exact shift. \code{tune.iload_n} is gone
    (\code{BUCK_ILOAD_N1}, \code{BUCK_ILOAD_N2}, rebuild to change).
    Ordinary pass 157 $\to$ 167 core clocks, telemetry 213, tunables
    227--255 plus one VDIV. The deadline check and the CMP1 write are now
    one asm block, since GCC's register allocation had put a reload of
    the timer address between them. Three agents verified it (every path
    walked statically; the linked code emulated over 270\,389 passes and
    1801 arms, A2 bit-exact, fault 5 never on its own and on every forced
    late pass; an adversarial review): no firmware error. Their other
    findings are in Section~\ref{sec:followup28}. (The sizes were changed
    to 128 $\times$ 128 on 28 September, Section~\ref{sec:followup28}.) \\
\bottomrule
\end{xltabular}}

% ============================================================
\section{Follow-up Requests, 28 September}
\label{sec:followup28}
% ============================================================

{\small
\renewcommand{\arraystretch}{1.25}
\begin{xltabular}{\linewidth}{@{}>{\raggedright\arraybackslash}p{0.27\linewidth}L@{}}
\caption{Requests and verification findings, 28 September.}
\label{tab:followup28}\\
\toprule
\textbf{Request or finding} & \textbf{Answer and result} \\
\midrule
\endfirsthead
\toprule
\textbf{Request or finding} & \textbf{Answer and result} \\
\midrule
\endhead
Make the load estimate faster as $S$ increases: from the 64-cycle
average down to a 1-cycle average.
  & Your choices: the two stages while $S\le S_0$; above it one average
    over $N=64\gg L$ cycles, $L=\lceil\log_2(S/S_0)\rceil$ limited to 6
    ($N$ halves each time $S$ doubles: 32 up to $2S_0$, \dots, 1 above
    $32S_0$); $S_0=0.002$, live as \code{tune.iload_s0}. The level follows
    a rise of $S$ at once and falls by one only after $N$ cycles at the
    current $N$, so the average never reaches back across a step when $S$
    peaks some cycles into it (corrected below: not when $S$ peaks in
    its first cycle). $L$
    comes exactly from the float bits of $S$ and $S_0$ (no division), and
    one counter $h$ holds the level ($h=\min(h, 63\gg L_S)+1$,
    $L=\mathrm{clz}(h)-25$). Cost: 7 clocks per quiet pass (one compare),
    about 35 more while it runs. On the telemetry and tunables passes it
    would have pushed those passes to about 283 / 294 clocks (plus the
    VDIV wait); you chose to skip the update there: A2 keeps its value for
    that cycle and the level holds. Simulated with $\pm2$ counts: steady
    state unchanged bit for bit; a 2\,A step reaches 90\,\% on A2 after
    35\,$\mu$s instead of 102. New Live Expression
    \code{out.i_load_lvl} (0: the two stages). (Superseded on 28 September
    by the ladder from 128, last row of this table.) \\
Verifier: \texttt{make check}'s \code{it cc} rule accepted broken builds
(trip removed, condition inverted, duty or raw counter compared).
  & Fixed: the checker follows the data from the counter read through
    the subtraction to the compare, and follows the \code{bcs} to the
    fault-5 trip; it tests its own rule at every run. All four broken
    builds now fail it. \\
Verifier: every stop overran the ADCs (the 2 $\times$ 64-entry ring
reset, about 570 clocks, ran in the stop pass) and added 1--2 to
\code{adc_miss}.
  & Your choice: the rings are cleared in 8 parts, one per stopped cycle
    (13\,$\mu$s), and a start waits until they are empty. With the
    256-entry rings of 128 $\times$ 128 (below): 32 parts of the same
    8 + 8 words (53\,$\mu$s). \\
Verifier (all versions): a RUN within about two periods of a STOP, or a
debugger halt at the wrong moment, could arm in the next period with its
sample already waiting; the first pass then tripped fault 4.
  & Your choice: \code{arm()} also requires that ADC1 has not converted
    again, checked after the counter read (no race; about 3 clocks on the
    arming path). \\
Verifier: \texttt{make test} re-implemented the firmware's per-cycle
load-current sequence instead of running it.
  & The sequence is one function, \code{buck_iload_cycle()}, which
    \code{run_pass()} and the closed-loop test both call (same machine
    code). Mutants of it now fail the tests, as does a broken ring reset
    (the empty-window test now resets a full window). \\
Verifier: wrong or stale text.
  & Corrected: the $\hat{\bar{\imath}}_L$ noise factor, ``the last $N$
    cycles'', the meaning of \code{loops}, the \code{busy_max} guidance,
    the quiet-pass count, the sum limit, the counting rule of the clock
    table (a register-offset store is 2 clocks in the TRM, 1 in the
    table), and the limit of \code{trip_cnt} (a pass more than a period
    and 64 clocks late shows as fault 4). \\
Verifier: the first pass after an arm at the very end of the window has
little margin with 45 extra wait states per pass.
  & Refuted by a second agent: the case needs wait states that the
    board's own \code{busy_max} reading rules out, and a late arm cannot
    happen on its own. Unchanged. \\
Make the load-current average 128 $\times$ 128; does it cost CPU?
  & No extra instructions: the rings hold running totals, so every sum is
    one subtraction whatever $N$. The linked loop has the same instruction
    sequence on every path as with 32 $\times$ 32, with equal cycle counts
    (constants, offsets and some encoding widths differ, plus two padding
    \code{nop}s that never run; same 4658 bytes; \texttt{make check}
    passes). The encodings moved the code by $\pm2$ bytes: the first
    instruction of every steady-state (HOLD) pass no longer starts on a
    word boundary, which may cost up to 1 clock (6\,ns) per pass; you chose
    to accept and document it. RAM: the rings grow from 0.5 to 2\,KB (91\,176\,B of
    96\,KB). Your choices (the ladder was moved below $S_0$ later, last
    row): the faster averages start at 128
    ($N=256\gg L$, $L$ limited to 8: 128 up to $2S_0$, \dots, 1 above
    $128S_0$; the two stages resume 255 counter steps after a 1-cycle
    average, 263--264 passes with the held telemetry and tunables passes),
    and the stop clears the rings in 32 parts of 8 + 8 words. A load step
    now takes 255 cycles (425\,$\mu$s) through the two stages. The sums
    stay exact for a mean below 20\,A (320\,A before); above it A2
    wraps. With the default $\alpha_I$ that is out of reach; with
    $\alpha_I$ of about 0.15 or more and an A1 reading the current loop
    cannot pull down it can happen (you chose to document it). Simulated
    with $\pm2$ counts: 0.0014 / 0.0010\,A rms at 1 / 3\,A (0.0023 /
    0.0015 at 32 $\times$ 32), 1.6 times less for a 3.4 times slower step:
    the A0 part falls by 1448 from the per-cycle value, but the correlated
    $\hat{\bar{\imath}}_L$ part, which dominates at both sizes, only by
    about 1.9 from 32 $\times$ 32. A 2\,A step reaches 50 / 90\,\% on A2
    after 18 / 40\,$\mu$s (5 / 35 before; highest level $N=8$, 2 before;
    superseded by the last row).
    The unused RAM-execution linker script no longer fits (40\,B over
    already before this change; now 1576\,B). \\
Verification of the 128 $\times$ 128 change (five agents, each finding
checked by a second).
  & No firmware error: every changed constant in the linked code is the
    value the 128 $\times$ 128 source calls for (so the earlier emulation
    of the 32 $\times$ 32 build carries over), the level counter was
    checked for every $h$ and $t$, the stop clear writes every slot
    once. Found and corrected: the alignment clock above; test gaps (the
    window test could overflow after a length change on other random
    seeds, the noise check did not scale with $N$ and failed at 32
    $\times$ 32 and below, the step bound failed at 16 $\times$ 16 and
    below, nothing tied the ladder to 128, the stop-clear test missed
    the last ring slot); and doc errors (the return time, the
    frequency below which content passes, which noise part dominates,
    and older rows not marked as superseded). \\
First board capture of the 128 $\times$ 128 build: why does A2 dip after
rising on a 1 $\to$ 3\,A step?
  & In simulation the dip appears when $S$ peaks in the step's first
    cycle: a $v_\mathrm{out}$ jump across the output capacitor's ESR
    makes the first-cycle $\Delta v$ large. The level then drops back
    ($N$ doubles every $N$ cycles), and each doubled average reaches back
    across the step (about 1\,A with 5\,m$\Omega$ ESR); with the ladder
    from 128 the last such window slides past the step only after about
    0.2\,ms. A real capacitance different from the model's 270\,$\mu$F
    also bends the per-cycle estimate itself during the $v_\mathrm{out}$
    transient. To tell the two apart on the board: \code{tune.iload_s0 =
    0} shows the per-cycle estimate on A2 (use scope averaging triggered
    on the load step), \code{tune.iload_s0 = 1e30} the two stages
    alone. \\
Board: with $S_0 = 100$ a smooth rise, with $S_0 = 0.01$ the dip. What
causes it; would capping $S/S_0$ help?
  & $S_0$ does not change the estimate, only which average A2 shows, so
    the dip comes from the ladder: with $S$ only a few times $S_0$ its
    shortest window is long (64 cycles at $S/S_0\approx 2$), and when it
    doubles it reaches back across the step (simulated with a clean
    estimate, the model plant: 0.72\,A at $S_0=0.01$, 0.05\,A at 0.002,
    none at 100).
    Capping $S/S_0$ keeps the windows long and does not help; a floor on
    the window does (8 or 16 cycles: 0.07--0.09\,A at every $S_0$). \\
Design: normally 128 $\times$ 128; as $S$ approaches $S_0$ the window
drops to 1.
  & Your choices: the window shrinks below $S_0$ and reaches 1 cycle
    once $S$ passes $S_0$ (the ladder moved down by 128: the two stages
    up to $S_q=S_0/128$, 128 cycles up to $S_0/64$, halving each time $S$
    doubles, 2 up to $S_0$, 1 above), and $S_0 = 0.128$ (0.002 before),
    so that steady state stays on the two stages (simulated $S\le 0.001$
    there). No extra instructions: $S_q$ is folded into the one saturating
    subtraction, which made the level update 1 clock shorter with the same
    instructions (a 1-clock immediate move replaced a 2-clock literal
    load), and the gate compares
    $\mathrm{bits}(S) < \mathrm{bits}(S_q)+1$. Simulated with $\pm2$
    counts: steady state on the two stages in every scenario (1.4 /
    1.0\,mA rms); a 2\,A step goes to a 4-cycle average and reaches 50 /
    90\,\% after 12 / 35\,$\mu$s. Still possible: a dip when $S$ peaks only
    just above $S_q$ (a 0.5\,A step falls back by about 0.3\,A, as
    before) or in the step's first cycle, and an estimate that itself
    overshoots and falls back if the real capacitance or ESR differ from
    the model (0.55\,A simulated with 10\,m$\Omega$ and twice the
    capacitance, with either ladder). This is the earlier ladder with its
    start at $S_0/128$: $S_0 = 0.256$ is exactly the earlier 0.002. The
    margin is small: $S_q = 0.001$ is only 2--3\,\% above the simulated
    peak steady-state $S$ (up to 0.0012 with other noise seeds). On the
    board, \code{out.i_load_lvl} should read 0 in steady state; if not,
    raise $S_0$. Verified by four agents, each finding checked by a
    second: no firmware error; the fast passes are 1 clock shorter; the
    gate and ladder match your thresholds for all $2^{32}$ values of $S$;
    a test now drives the gate through \code{buck_iload_cycle()}. (A floor of 1 cycle above $S_0$ was built first and
    replaced by this at your request.) \\
\bottomrule
\end{xltabular}}

% ============================================================
\section{Follow-up Requests, 24--25 October (BUCK\_Adpt\_Learning)}
\label{sec:followup1024}
% ============================================================

{\small
\renewcommand{\arraystretch}{1.25}
\begin{xltabular}{\linewidth}{@{}>{\raggedright\arraybackslash}p{0.27\linewidth}L@{}}
\caption{Requests, 24--25 October.}
\label{tab:followup1024}\\
\toprule
\textbf{Request} & \textbf{Answer and result} \\
\midrule
\endfirsthead
\toprule
\textbf{Request} & \textbf{Answer and result} \\
\midrule
\endhead
How is $i_\mathrm{avg}$ obtained; does it converge to the measured
average?
  & $i_\mathrm{avg}$ (\code{out.i_avg}) is the estimator's corrected
    value: the Pad\'e model predicts the cycle average from the corrected
    valley current, the duty, $V_\mathrm{in}$ and $v_\mathrm{out}$, and the
    measured average (four ADC2 samples per cycle, summed in hardware)
    corrects it by $\alpha_I = 0.0328$. It settles, but to the measured
    value only if the model matches the plant: with a model bias about
    93\,\% of it remains ($A_I\approx0.56$ per cycle, proportional
    correction only). Simulated with body-diode dead time: measurement
    within $\pm0.03$\,A, estimate $-1.36$ / $+0.34$\,A off. Not changed. \\
Copy the code and structure to the new project BUCK\_Adpt\_Learning.
  & Copied (sources, CMSIS, tests, tools, docs, the custom linker
    scripts); the Makefile builds \code{BUCK_Adpt_Learning.elf}; after
    CubeIDE was closed its \code{.cproject}, \code{.settings} and a
    renamed debug launch were installed (the original's, with the project
    name replaced). \\
From the PLECS model: Pad\'e as estimator + correction; the load
low-pass from averaging to a first-order term damped by the score;
score with $\rho = 0$ and $\Delta d$ instead of $\Delta i_L$.
  & The PLECS model already has all three (score
    $(e/0.02)^2+(\Delta d/2)^2$, $\rho=0$, $S_\mathrm{th}=4$; Adaptive Low
    Pass with $\tau$ from 33.32\,$\mu$s to 0 at $S\ge S_\mathrm{th}/2$; a
    3-state predictor/corrector). Your choices: port to the firmware copy;
    no estimator change now (the firmware's Pad\'e corrects $i_L$ only);
    the filter with $\alpha$ interpolated (no division per pass):
    $\alpha = \alpha_0 + (1-\alpha_0)\min(\max(S/(0.5S_\mathrm{th}),0),1)$,
    $\alpha_0 = T/(\tau_0+T)$, $\tau_0$ live as \code{tune.iload_tau0}
    (33.32\,$\mu$s); the firmware's scales ($e_v/1$\,V, $\Delta d/1$,
    $S_\mathrm{th} = 1$). The 128 $\times$ 128 average, its ladder, the
    rings and the stop-time clearing are gone. The coefficients take one
    division, in the stopped cycles and in the telemetry pass (not the
    tunables pass, the longest). Static count (one model for both):
    normal pass 156 / 169 / 161 / 169 clocks against 167--209 before,
    telemetry 243--260 (219--247; 220--248 with the dependent-load rule
    added on 25 October), tunables 232--264 (240--268); each 28 September
    figure is 1 clock more if a load right after a register-offset store
    is not counted as pipelined, as the TRM has it (open); RAM
    $-2$\,KB. Simulated with $\pm2$ counts: A2 noise 6.6 / 6.3\,mA rms
    (1.4 / 1.0 with 128 $\times$ 128); a 2\,A step reaches 50 / 90\,\%
    after 27 / 60\,$\mu$s (35 / 103 with $\alpha$ fixed); with
    $S_\mathrm{th} = 1$ the step's score (about 0.04) lifts $\alpha$ only
    to 0.13. Host tests rewritten for the filter and $\Delta d$; all
    pass. Not run on the board. \\
Verification of the 24 October changes (four agents, each finding
checked by a second).
  & The filter arithmetic holds for every input (coefficients finite for
    all $\tau_0$, $S_\mathrm{th}$; $\alpha$ in its range for every $S$),
    and the clock figures were reproduced. Fixed: the coefficients were
    zero, and $\alpha$ therefore 0, for the first 24--56 passes when
    \code{command} was set before the warm-up ended (e.g.\ at the
    \code{main} breakpoint), and a halt's tunables were skipped at the
    start: they are now formed at boot and in every stopped cycle that
    does not start the run. Two test gaps (the score's $d$ against the
    CMP1 duty; the duty at arming computed independently) and several
    comments and figures. Left for you: A2 starts near 0\,A after arming
    (the gates were off) and rises with $\tau_0$, about 100\,$\mu$s;
    whether PLECS's $\Delta d$ at a rise is this update's or the
    previous one (not confirmed); the README lists the other differences
    to PLECS. \\
How to set $\tau$?
  & Write \code{g_buck.tune.iload_tau0} in seconds (Live Expressions
    row; default 33.32e-6). While running it takes effect within about
    96 cycles (160\,$\mu$s); 0 means no filtering ($\alpha = 1$), values
    above 1\,s count as 1\,s. \\
An additional averaging filter on the load estimate, 32 samples; then
``remove the 32 avg''.
  & Built as asked (before the first-order filter, always on: a ring of
    the last 32 estimates and their exact running sum). Measured: 12--13
    clocks in every pass, the longest pass 276 of 280 clocks (280 with
    the loop head); simulated A2 noise 2.1 / 1.4\,mA rms instead of 6.6 /
    6.3, a 2\,A step at 90\,\% after 98 instead of 60\,$\mu$s. Removed at
    your request before it was verified or run: the code, tests and
    README are back to the filter alone; the rebuilt control loop is
    identical, instruction for instruction, to the build before the
    average (4482\,B), and the host tests pass. \\
25 October: ``change the score function to $a u^2 + b u^2 + c$, or any
other suggestion; I need it to increase quadratically in a transient'',
and a minimum of 3.3\,$\mu$s for the low-pass filter coefficient.
  & Your decisions: $S = a e_v^2 + b e_v + c$, $e_v = v_\mathrm{out} -
    V_\mathrm{ref}$ (signed, as before); $a$, $b$, $c$ live as
    \code{tune.score_a}, \code{tune.score_b}, \code{tune.score_c},
    defaults 1, 0, 0 ($S = e_v^2$, bit for bit); $\rho$, \code{ev_norm},
    \code{du_norm} and the $\Delta d$ term removed, with the previous duty
    (\code{st.d}) that only $\Delta d$ used. The one $S$ drives D4
    ($S \ge S_\mathrm{th}$) and the filter. It is evaluated as
    $(a e_v + b) e_v + c$: two multiplies, two adds, no division, within
    4 roundings of the formula. With $a < 0$, $b \ne 0$ or $c < 0$, $S$ can be
    negative: then $\alpha = \alpha_0$, and D4 stays low while
    $S_\mathrm{th} > S$; a NaN $S$ gives $\alpha_0$ and D4 low, $+\infty$
    gives $\alpha = 1$ and D4 high. Static count, HOLD / FREE\_UP / UP /
    DOWN, with a load whose address comes from the load just before it
    counted as not pipelined (Cortex-M4 TRM; the HOLD and FREE\_UP copies
    each have two such loads, the 24 October loop none): normal pass
    150 / 154 / 160 / 160 clocks (156 / 161 / 169 / 169), telemetry
    239--253 (243--260), tunables 211--243 (232--264: the VDIV is gone),
    the CMP1 write 74 / 74 / 76 / 76 clocks into the pass (77); the first
    pass after arming 310--336 of its 416 clocks (334--357 as counted on
    24 October, 334--356 by the same count), and
    \code{BUCK_ARM_LATEST} stays at 4096 ticks (about 80 clocks of margin,
    60 before). Counting those loads as pipelined, as first reported,
    gave 149 / 153 / 160 / 160, 237--253, 210--243, 73 / 73 / 76 / 76
    and 309--336. With the TRM rule two other orders of the four copies
    are no worse in any pass length or the CMP1 write and better in most
    (README); the order is unchanged until you decide. With the default
    coefficients every simulated A2 figure is unchanged; only the steady-state peak $S$ changed (no
    $\Delta d$ term). New tests: the formula against double over 2 million
    random $a, b, c, e_v$; NaN, $\pm\infty$ and negative scores; closed
    loops with $a = 100$, and with $a, b, c = 100, 5, -0.25$ ($S = -0.25$
    at rest, D4 only on the overshoot of the step down); with the default
    score and filter, $\alpha$ below 0.1 on average and A2 noise below
    10\,mA rms. \texttt{make check} now also requires every tunable to be
    read at the loop entry and in the tunables pass (it fails when the
    tunables pass stops copying $a$, $b$, $c$, or the telemetry pass stops
    forming the filter coefficients); the rest of the wiring in
    \texttt{buck.c} (e.g.\ which value feeds the filter as $S$, or which
    value drives D4) is checked by review only. \\
``Don't cap it; use different $S_\mathrm{th}$ for score and filtering.''
  & The 3.3\,$\mu$s minimum $\tau$ was considered (it would have capped
    $\alpha$ at $T/(3.3\,\mu\mathrm{s}+T) = 0.336$) and dropped: $\alpha$
    still reaches 1 at $S \ge S_f$ (at $S = S_f$ exactly, float rounding
    can leave it up to $2.4\times10^{-7}$ below 1). The filter has its own
    live threshold \code{tune.iload_s_th}, $S_f = 0.5\,$\code{iload_s_th}, default 1 (the
    $S_\mathrm{th}$ default, so A2 is unchanged until it is tuned), with the
    clamping and NaN rules $S_\mathrm{th}$ had there; \code{tune.s_th} sets
    D4 only. Tests: D4 and the filter each follow their own threshold, bit
    for bit unchanged when the other one moves; closed loop with
    \code{iload_s_th} $= 0.01$: $\alpha = 1$ on the steps, A2 at 90\,\%
    after 8\,$\mu$s, D4 low. Live Expressions: \code{tune.rho},
    \code{tune.ev_norm}, \code{tune.du_norm} replaced by
    \code{tune.score_a}, \code{_b}, \code{_c} and \code{tune.iload_s_th}
    (42 rows). Not run on the board. \\
\bottomrule
\end{xltabular}}

% ============================================================
\section{Follow-up Requests, 25 October (Buck\_disturb\_obs)}
\label{sec:followup1025obs}
% ============================================================

You asked for your fixed-gain load-current observer, studied it against
the score-damped filter of BUCK\_Adpt\_Learning, and decided how it goes
into a new copy, Buck\_disturb\_obs. The study's numbers come from a
closed-loop host simulation (the same plant as \texttt{make test}, the
firmware's own per-cycle code); those marked ``this build'' are from the
code as delivered. Nothing has run on the board.

{\small
\renewcommand{\arraystretch}{1.25}
\begin{xltabular}{\linewidth}{@{}>{\raggedright\arraybackslash}p{0.27\linewidth}L@{}}
\caption{Requests and decisions, 25 October (Buck\_disturb\_obs).}
\label{tab:followup1025obs}\\
\toprule
\textbf{Request or decision} & \textbf{Answer and result} \\
\midrule
\endfirsthead
\toprule
\textbf{Request or decision} & \textbf{Answer and result} \\
\midrule
\endhead
Your observer: $e = v_\mathrm{out} - \hat v$,
$\hat v \leftarrow \hat v + a(i_L - \hat I) + \beta e$,
$\hat I \leftarrow \hat I - \gamma e$; compare it with today's filter.
  & Study (two simulations, each rebuilt by a second agent). In
    predict-then-correct order (PC) the observer is exactly two identical
    first-order stages $y \leftarrow y + (1-p)(\hat I_k - y)$ on the
    per-cycle estimate $\hat I_k = i_L - (C/T)(v_{k+1} - v_k)$, transfer
    function $(1-p)^2/(1-pz^{-1})^2$, for $\beta = 1-p^2$ and
    $a\gamma = (1-p)^2$, $a = T/C$. It has the input's DC value, so fed
    $\hat{\bar\imath}_L$ it keeps the dead-time bias ($-1.36$ / $+0.34$\,A
    at 1 / 3\,A with body diodes); fed the measured average $i_\mathrm{meas}$
    the bias is $+0.027$ / $-0.018$\,A. At $\tau = 9.2\,\mu$s (fed
    $\hat{\bar\imath}_L$) it matched the filter's body-diode step (90\,\%
    after 60.0 / 73.3\,$\mu$s) with 5.09 instead of 6.40\,mA rms of A2
    noise; fed $i_\mathrm{meas}$ the noise is 14--15\,\% higher. A prototype
    with live gains was 8--13 clocks shorter per normal pass than the filter
    (on the code before the score change). \\
1.\ Base: BUCK\_Adpt\_Learning after the score change and the ``fix
base'' stage; the score keeps driving D4.
  & Copied as such. $S = a e_v^2 + b e_v + c$ (\code{tune.score_a},
    \code{_b}, \code{_c}, defaults 1, 0, 0) drives D4 ($g = S \ge
    S_\mathrm{th}$) and nothing else: with the filter gone it no longer
    touches A2. \\
2.\ The load estimate on A2 and \code{out.i_load} from your observer,
predict then correct, $a = T/C$, $v$ = this pass's two-sample mean, $i_L$
= the raw 4-sample ADC2 mean of the cycle, after the guarded CMP1 write;
A2 code as before.
  & Done: \code{buck_iload_obs()} in \code{run_pass()} after
    \code{buck_ctrl_slow_k()} (so after the CMP1 write), fed
    \code{st.v} and \code{st.im} of the pass's \code{buck_step_t};
    A2 $=$ \code{buck_iload_dac(}$\hat I$\code{)}.
    $a = $ \code{BUCK_OBS_A} $= (T/C)$ as float $= 1/162$\,V/A. It replaces
    the per-cycle $\hat I$ (\code{buck_load_estimate}) and the filter; the
    controller state's previous $v_\mathrm{out}$ (\code{buck_state_t.v},
    i.e.\ \code{g_buck.st.v}), used only by $\hat I$, went with it. \\
3.\ Live gains \code{tune.obs_beta}, \code{tune.obs_gamma}; defaults
from a double pole at $\tau = 9.2\,\mu$s.
  & $p = e^{-T/\tau} = 0.83430$, $\beta = 1 - p^2 = 0.3039396$,
    $\gamma = (1-p)^2 C/T = 4.447826$\,A/V ($4.4478256$; the float is
    $4.4478254$; the request's 4.447825 is truncated), formed at compile time from
    $\tau$ (\code{BUCK_OBS_TAU_S}) with the 7-term series already used for
    $\alpha_I$ (relative error $1.7\times10^{-10}$); the floats are the same
    bits as with libm's \texttt{exp} (checked by \texttt{make test} from
    $\tau = 9.2\,\mu$s written in the test, and against the digits). \\
4.\ A NaN or out-of-range pair rejected where the gains are reloaded
(tunables pass, loop entry), previous values kept, the compiled defaults
at the entry; out of range $=$ outside the stable region, derived by Jury,
closed boundary rejected.
  & Error dynamics $x_1' = (1-\beta)(x_1 - a x_2)$,
    $x_2' = \gamma x_1 + (1 - a\gamma) x_2$ (the same matrix moves the
    observer's own state), $P(z) = z^2 - (2 - \beta - a\gamma) z + (1-\beta)$.
    Jury: $a\gamma > 0$, $4 - 2\beta - a\gamma > 0$, $|1-\beta| < 1$, i.e.\
    $\gamma > 0$ and $0 < \beta < 2 - a\gamma/2$. No further condition: the
    study's extra $1 - a\gamma > -1$ is not needed ($a\gamma = 3$,
    $\beta = 0.25$: both poles at $|z| = 0.866$). The pair is judged as a
    pair; at the entry a rejected pair gives the defaults, in the tunables
    pass the gains in use stay. The test is exact in float for every pair
    (\code{buck_obs_gains_ok}: the rounding error of $a\gamma/2$ by Dekker's
    product, no FMA); a plain float test would accept 32 of the 80 points
    exactly on the boundary that \texttt{make test} constructs. Cost (this
    build): 28--38 clocks in the tunables pass, 24--34 on the arming path. \\
5.\ No NaN guard on the observer states.
  & None. \\
6.\ Remove \code{tune.iload_tau0}, \code{tune.iload_s_th},
\code{out.i_load_alpha} and what served only the filter; adjust rule 4 of
\texttt{make check}.
  & Removed, with \code{buck_lpf_k_t}, \code{buck_iload_lpf_k},
    \code{buck_iload_alpha}, \code{buck_iload_lpf}, the VDIV in the
    telemetry pass and the stopped loop, the \code{BUCK_ILOAD_*} constants,
    the filter's tests and the closed-loop $\alpha$ bounds. Rule 4 now
    requires each tunable to be read at the loop entry \emph{and} in the
    passes, counted separately (until now: two reads anywhere). Reason: GCC
    repeats the gains' check in all four copies of the tunables pass, so an
    entry taking the defaults without reading the tunables still had four
    reads and passed; it now fails, as does a tunables pass without the
    gains' reload. It also requires \code{obs_beta} and \code{obs_gamma}
    among the fields, and (repair after the verification) that they are not
    read on the path of an ordinary pass (from a pass head, without the
    branch of the \texttt{loops \& 31} test), so the reload line moved into
    every pass fails it too. A check run every pass on the loop's own copy
    \code{c->t} reads nothing from \code{g_buck} there and passes (review
    only). \\
7.\ Seeding: at the loop entry $\hat I$ = the arming pass's
$i_\mathrm{meas}$, $\hat v$ = its $v$.
  & \code{buck_obs_seed(&b->out)} at the entry (arming publishes
    \code{out.i_load} $= i_\mathrm{meas}$ as before). \\
8.\ Project setup: copy BUCK\_Adpt\_Learning's code and structure into the
new empty project as on 24 October; prepare \code{.cproject},
\code{.settings} and the debug launch for installation after CubeIDE is
closed.
  & Copied: \code{Inc}, \code{Src} (the template's \code{main.c} replaced;
    \code{syscalls.c}, \code{sysmem.c} and the startup file were identical),
    \code{Startup}, \code{Drivers}, \code{Tests}, \code{tools}, \code{docs},
    \code{README.md}, the Makefile (\code{NAME := Buck_disturb_obs}),
    \code{.gitignore} and the custom linker scripts (over the template's:
    RAM 96\,K instead of 128\,K, and the \code{.ccmtext} output section).
    The project's \code{.project} kept; \code{.cproject} and \code{.settings}
    untouched. Prepared outside the project: BUCK\_Adpt\_Learning's
    \code{.cproject}, \code{.settings} and debug launch with the project name
    replaced (launch renamed), with diffs against the template. Not yet
    installed. Before any change the copy built to the same binary as
    BUCK\_Adpt\_Learning. \\
9.\ BUCK\_Adpt\_Learning stays as it is; the projects are independent.
  & Not modified by this stage. \\
Results (this build).
  & \texttt{make}: 0 warnings, RAM 89\,112\,B, FLASH 8504\,B;
    \texttt{make check} ok; \texttt{make test} 0 failures. Static count
    (HOLD / FREE\_UP / UP / DOWN; a load whose address comes from the load
    before it as 2 clocks), against BUCK\_Adpt\_Learning: normal pass 140 /
    142 / 148 / 148 (150 / 154 / 160 / 160); telemetry 180 / 182 / 190 /
    188 (up to 253: no VDIV now); tunables 237--267 (211--243: the gains'
    check); CMP1 write 73--75 clocks into the pass (74--76); first pass
    after arming 337--368 of 416 (310--336). Longest pass 267 of 280 (the
    tunables pass). No order of the four copies is at least as good on
    every measure; UP, FREE\_UP, DOWN, HOLD is better on the limiting
    figures (longest pass 265, arming 336--367, the CMP1 write 1 clock
    earlier in every copy) and 2 clocks longer on DOWN's telemetry pass
    (190 against 188); DOWN, FREE\_UP, UP, HOLD likewise, with arming
    335--366 and UP's tunables pass 265 (262). Averaged over 64 passes in steady state (HOLD): 142.3
    clocks (152.4). Simulated with $\pm2$ counts of ADC noise: A2 noise 5.7
    / 5.7\,mA rms at 1 / 3\,A (bound 8\,mA), bias $+0.024$ / $-0.020$\,A
    ($+0.027$ / $-0.018$ without noise, $-0.015$ / $-0.017$ with ideal dead
    time; bound 0.05\,A); a 2\,A step at 50 / 90\,\% after 15.0 / 88.3\,$\mu$s
    up and 13.3 / 41.7\,$\mu$s down (body diodes); the observer within
    $5\times10^{-6}$\,A of the two first-order stages. On the study's
    2\,kHz 1--3\,A square profile (its harness, not a host test) A2's rms
    error is 380\,mA with body diodes and with ideal dead time, against 1087
    and 445\,mA for BUCK\_Adpt\_Learning's filter. Mutants: the observer
    in the other order, $\gamma$'s sign, the check without its rounding
    error, accepting the boundary, without $\gamma > 0$ or $\beta > 0$, the
    reload not keeping the previous gains, a wrong seed, default $\gamma$ or
    $a$, $\tau$ = 9.3 or 10\,$\mu$s: each fails \texttt{make test}; the
    reload or the entry check removed from \code{buck.c}, or the reload
    line moved into every pass: each fails \texttt{make check}. Still
    review only (each passes \texttt{make}, \texttt{make check} and
    \texttt{make test}), for example: the check run every pass on
    \code{c->t}, or in the telemetry pass; the observer stepped only in the
    tunables pass, or twice per pass; $\hat I$ seeded from
    \code{out.i_avg}; the observer fed $\hat{\bar\imath}_L$ or $V_\mathrm{ref}$,
    or moved before the CMP1 write; the check's result ignored; $\hat v$
    started at 0 in \code{run_loop}; A2 written from $\hat v$;
    \code{out.i_load} not updated in the telemetry pass (an uninitialised
    value published, no warning) or published as $\hat v$ or
    $i_\mathrm{meas}$; the entry falling back to 0 instead of the defaults;
    the reload checking the previous tunables' pair. \\
Left for you.
  & Decided by you after the report (25 October), all kept as built: the
    exact check stays at the loop entry (arming margin 48 clocks instead of
    80); no telemetry of the gains in use after a rejected write; the copy
    order UP, FREE\_UP, HOLD, DOWN; the wiring mistakes above stay review
    only, without further \texttt{make check} rules. Still open: with ESR the model misses $162 \times$ ESR of the
    capacitor current (5\,m$\Omega$: the step down at 90\,\% after 130
    instead of 42\,$\mu$s). The 2\,kHz profile is not among the host
    tests. The 28 September column of the clock table counts a
    load right after a register-offset store (and, for arming, after an
    STRD) as pipelined, against the TRM: restate it (+1 clock per figure)
    or keep it with the note. Until the prepared settings are installed, a
    CubeIDE build of this project fails (no CMSIS include paths). \\
\bottomrule
\end{xltabular}}

% ============================================================
\section{Follow-up Requests, 25 October evening (Buck\_disturb\_obs: fault 5 on the board)}
\label{sec:followup1025f}
% ============================================================

The first board run of Buck\_disturb\_obs (CubeIDE Debug build, the same
CCM code as the Makefile build) tripped fault~5 on the first pass after
arming. You had the cause analysed, chose option~(f) ``setup before
arming'' and asked for timing stamps. Clock figures are counted on the
linked ELF from the issue of one instruction to the issue of another,
static (your model) and with the wait states of the boot benchmark
\code{g_buck.bench} (nominal, and pessimistic with CCM literals at
$+2$ and no decay of the write penalties). Nothing of the fix had run on
the board when this section was written; it ran on 27 October (you
reported that it works; no data were read). Its clock figures are
superseded by the 28 October build (Section~\ref{sec:followup1028}).

{\small
\renewcommand{\arraystretch}{1.25}
\begin{xltabular}{\linewidth}{@{}>{\raggedright\arraybackslash}p{0.27\linewidth}L@{}}
\caption{Requests and decisions, 25 October evening (the fault~5 fix).}
\label{tab:followup1025f}\\
\toprule
\textbf{Request or decision} & \textbf{Answer and result} \\
\midrule
\endfirsthead
\toprule
\textbf{Request or decision} & \textbf{Answer and result} \\
\midrule
\endhead
Board result (25 October, about 21:12).
  & Fault~5 on the first pass after arming: \code{loops} $= 8$ (its start
    value), \code{trip_cnt} $= 8608$ ticks against the deadline 8448
    (5~clocks late), $v_\mathrm{out}$ 0\,V, \code{adc_miss}, \code{warmup}
    and \code{busy_max} 0. \code{bench} (extra clocks per access over
    SRAM1): CCM literal or data load while executing from CCM $+1.06$;
    HRTIM counter read $+3$; ADC12\_CSR, ADC1\_DR, ADC2\_DR reads $+1$;
    DWT\_CYCCNT 0; a peripheral write followed at once by an SRAM or
    peripheral load $+2$ (by a CCM literal $+1$); back-to-back CMP1
    writes $+1.06$; DAC and BSRR about 0. \\
Analysis.
  & The static count had no wait states. With them the old path from
    \code{arm()}'s counter read to the first pass's guarded counter read
    was 334--365 static, 375--408 nominal, 405--440 pessimistic (the
    board's: 361 / 402 / 432; 30--32 CCM literal loads). The deadline is
    absolute (one period plus 264 clocks after the start of the arming
    period), and the trip means about 460--485 clocks at the usual arming
    point: 60--80 clocks (16--20\,\%) that the model does not explain.
    Hypothesis, unproven: VLDR/VSTR at 2~clocks, not pipelined ($+63$ on
    that path). Steady state keeps 100--150 clocks of CMP1 margin. \\
1.\ Fix (f), in Buck\_disturb\_obs only: \code{buck_run} calls
\code{run_loop} in the stopped cycle that sees RUN; the entry runs
there; then the loop waits for the next sample and arms with today's
conditions; a refused arm waits; a bad sample or a STOP leaves as
today; every stopped-loop duty that still applies is kept; arming with
the tunables copied at the entry is accepted.
  & Done. \code{buck_run} calls \code{run_loop(b)} where it called
    \code{arm()} (no bad sample, \code{warmup} 0, \code{state} STOPPED,
    \code{command} RUN). The entry (prologue, the exact gains' check as its
    first statement, \code{c.t}, Ki$\cdot$T, score coefficients, target,
    step, ADC base, \code{loops}) ends with an empty \code{asm volatile}
    taking those values as inputs, with a memory clobber: in the ELF the
    entry ends at 0x100004e2, before the wait's first status read at
    0x100004e6, with every tunable read and Ki$\cdot$T, the target and
    \code{loops} formed; nothing of it after the wait (without the
    \code{asm} GCC moves Ki$\cdot$T, the target and \code{loops} there:
    17~clocks more). Per waiting sample, as the stopped loop: the status
    wait with its 840-poll timeout (fault~4), the two DR reads, the raw
    codes published; a bad sample clears the overrun flags, counts in
    \code{adc_miss} and returns (still stopped); a STOP returns. The
    warm-up count and the state test do not apply (the call requires
    \code{warmup} 0 and STOPPED, and nothing changes them before the
    arming). Arming: CYCCNT and the counter read in one \code{asm}, the
    window 897--4095, then the ADC1 re-check (read after the counter);
    \code{buck_ctrl_arm} on the loop's \code{c.t} and \code{c.s};
    CMP1, then OENR; \code{state} $=1$; \code{out.*} published; the
    observer seeded from the arming sample; then the first pass.
    \code{g_buck.st} is written back at a stop or trip as before, no
    longer at arming (during a run it holds the previous stop's state).
    BUCK\_Adpt\_Learning and Buck\_Current\_Mode untouched. \\
2.\ Keep: the exact check literally at the entry and in the tunables
pass (defaults at the entry); the seeding; predict then correct;
$i_\mathrm{meas}$; $\tau = 9.2\,\mu$s; no NaN guard; the score and D4;
the copy order UP, FREE\_UP, HOLD, DOWN; \code{BUCK_ARM_LATEST} 4096;
fault~5 for any late CMP1 write, the first pass's included.
  & All kept; no line of code in \code{Inc/buck_ctrl.h} changed (two
    comments), and \texttt{make test} prints the same output, line for
    line, as the build before the fix. The first pass's CMP1 write has the
    same guarded check and trip as every other. \\
3.\ Timing stamps appended after \code{trip_cnt} (existing offsets
unchanged); per-pass code of the ordinary passes unchanged; stamp work
off the critical path; how to read them; Live Expressions rows checked
with gdb.
  & \code{arm_cnt}, \code{arm_cyc}, \code{first_cnt}, \code{first_cyc},
    \code{trip_cyc} at offsets 228--244 (\code{g_buck} 248\,B, at
    0x20000000 as before). Each CYCCNT is read by the instruction just
    before its counter read (both pairs in one \code{asm}), so
    $\mathtt{first\_cyc} - \mathtt{arm\_cyc}$ equals
    $(\mathtt{first\_cnt} + n \cdot 8960 - \mathtt{arm\_cnt})/32$ within a
    clock or so, $n$ being the period starts between the two reads: $n = 1$
    when the check falls in the period after arming, as it normally does.
    Only then is $8448 - \mathtt{first\_cnt}$ the first write's margin;
    \code{first_cnt} below 2048 means the check came more than a period
    late (fault~5), and that row then reads 6401--8448 (8448 also before
    the first start, with \code{first_cnt} 0). The first pass is a
    fifth copy of the pass (its reference step chosen at run time); it
    alone reads the stamp CYCCNT and stores the four stamps after its CMP1
    write (on its trip or on a stop before its check too;
    \code{0xFFFFFFFF} in \code{first_*} for the latter, where
    $\mathtt{first\_cyc} - \mathtt{arm\_cyc}$ means nothing), but not when
    it ends in \code{buck_fatal()} (fault~4 by the sample timeout, or
    fault~8): the stamps then keep the previous start's values. 0 until the
    first start. \code{trip_cyc}: CYCCNT read at the start of the trip
    branch of every pass (its first instruction in the first pass; in the
    ordinary copies after two register moves and a stack load, and in three
    of them a branch), 11--16 clocks after the counter read in
    \code{trip_cnt}; 0 until a fault~5. The four ordinary copies contain
    no stamp code, but they are not instruction-identical to the build
    before (item~4). Live Expressions: seven rows
    added (the five stamps, \code{first_cyc-arm_cyc} and
    \code{8448-(int)first_cnt}), 48 in all, each evaluated with
    \texttt{arm-none-eabi-gdb} against the ELF. \\
4.\ Counts from the linked ELF; 280 clocks per pass; CMP1 by 264
clocks into the period.
  & Arming check to the first pass's guarded counter read (first pass
    starting late): 172--195 static, 196--221 nominal, 213--240
    pessimistic (before 334--365 / 375--408 / 405--440); the board's path
    195 / 221 / 240 (before 361 / 402 / 432), estimated 231--266 on the
    board (VLDR/VSTR at 2~clocks: 236; flat 16--20\,\%: 257--265; up to
    about 306 if the old excess, 56--85 clocks over the nominal count, is
    one fixed cost that falls on this path too). Against 416 at the latest
    arming point (4096 ticks) that leaves the arming path 176--220 clocks
    of slack (about 150 with the board's excess as a proportion, 110 or
    more as a fixed cost), which is not the first write's margin: the
    first pass's head comes 94--117 / 108--133 / 118--145 clocks after the
    arming check, at most 273 clocks into the arming period, before the
    next end of conversion (about 313), so at every accepted arming point,
    4096 included, it waits for its sample and checks at 120--138 clocks
    into the next period: 126--144 clocks of margin (about 108--130 with
    the board's excess). The arming CMP1 write (no deadline check by
    design, so no fault~5) is the tighter limit on
    \code{BUCK_ARM_LATEST}: 77--88 / 91--103 / 99--112 clocks after the
    arming check's counter read (\code{arm()}: 83--90 / 99--107 /
    109--118), so at 4096 it lands 24--37 clocks before 264 (pessimistic;
    about 13--31 with the board's excess as a proportion, past 264 if that
    excess were a fixed cost there); a late arming write leaves the first
    enabled period at the previous CMP1 (\code{BUCK_CMP_MIN}). GCC stores
    10 of the 11 \code{out.*} fields before that write (\code{arm()}
    stored 10 of them after OENR and \code{out.g} between CMP1 and OENR;
    with a memory barrier after OENR the write comes 6--7 clocks earlier,
    counted, not kept). The stamps' CYCCNT
    reads: the arming one moves the arming counter read 6--10 clocks later
    after its sample, not the write relative to that read; the first
    pass's costs 1~clock on its \code{strcc}. The first pass as a whole
    (body to the next copy's status read, that read included, as for the
    passes below): 193--217 static, 221--246 nominal, 240--266
    pessimistic; to HOLD 193--205, FREE\_UP 202--214, UP 201--217, DOWN
    195--207 static (191--215 / 218--243 / 237--263 without that read).
    It is the longest pass type after the tunables pass; with its poll
    head about 250 nominal, so with the board's 16--20\,\% it can run up
    to about 20 clocks past 280. That only delays the second pass (by up
    to about 20 clocks after its sample), whose CMP1 write still lands
    about 90--120 clocks before 264; nothing accumulates (a normal pass is
    about 150). The stopped cycle
    that sees RUN reaches the arming wait at most 210 clocks into the
    period (pessimistic), before the next end of conversion (about 313).
    Ordinary passes, static (before $\to$ now, HOLD / FREE\_UP / UP /
    DOWN): normal 140 / 142 / 148 / 148 $\to$ 140 / 143 / 147 / 147;
    telemetry 180 / 182 / 190 / 188 $\to$ 182 / 185 / 189 / 189;
    tunables at their longest 262 / 264 / 262 / 267 $\to$ 262 / 259 /
    267 / 264; CMP1 74 / 73 / 75 / 75 $\to$ 72 / 72 / 73 / 73;
    \code{busy_max} window 141 / 143 / 151 / 149 $\to$ 143 / 146 / 150 /
    150. Longest pass 267 (294 nominal, 309 pessimistic with the head),
    as before. The ordinary copies' normal passes keep their instruction
    count (131 / 133 / 138 / 138, head to loop-back) with other registers
    and stack slots, but are not the same instructions:
    \texttt{lsls~\#27} for \texttt{tst~\#31} in all four; in FREE\_UP
    the ADC2 code is the value spilled (the ADC1 code before) and is
    reloaded with \texttt{vldr} where \texttt{vmov} moved it, before the
    CMP1 write, a value held in \texttt{s31} before is reloaded from the
    stack after it, and one 16-bit \texttt{ldr} is 32-bit; instructions in
    another place: HOLD 7, UP 3, DOWN 3, FREE\_UP 11 (longest common
    subsequence, registers and stack slots abstracted). FREE\_UP's CMP1
    write is 1~clock earlier and its pass 1~clock longer. \code{run_loop} 4308\,B $+$ \code{arm()} 432\,B $\to$
    5620\,B; RAM 89\,112\,B unchanged; FLASH 8504 $\to$ 9396\,B. \\
\texttt{make check}.
  & Kept every rule. The arming CMP1 write, now inside \code{run_loop},
    is exempt from the deadline guard, and only it: the one CMP1 write
    followed within three instructions by the OENR write through the same
    base, not on the path of an ordinary pass. A literal pool between the
    \code{strcc} and its \code{bcs} (\code{b}, then \code{bcs}: GCC does
    this in the first pass) is accepted, nothing else. The entry of the
    tunables rule still ends at the first ADC12\_CSR read, which is now the
    arming wait's, found by tracking the entry's register values (GCC
    reads it through \code{ip}); so ``read at the loop entry'' means read
    before the arming wait. Each change is self-tested. On 26 October 13
    of 21 broken versions of these changes still passed the self-tests:
    two that cannot behave differently on \code{run_loop}, and the join
    with the paths in the other order; a sub, an mvn, a movt or a literal
    load; a register write or a call not making a register unknown; the
    entry continued past a call that does not return; two arming writes;
    one on the path of an ordinary pass; a branch between CMP1 and OENR.
    Cases for the eleven were added; of five more broken versions, three
    (a register copy not tracked right) passed those too, and cases for
    them were added. Also on 26 October: the exemption of stores to
    offset 156 through a register holding \code{g_buck}'s address (a
    \code{g_buck} field, not CMP1) no longer breaks when that register is
    written with a value never used as an address (every path from the
    write returns, branches out or calls \code{buck_fatal()} first). GCC
    used the register as a scratch on such a block when a bad sample
    while waiting is not counted; \texttt{make check} then reported its
    stores to \code{busy_max}, at CMP1's offset, as unguarded CMP1 writes
    and found no register holding \code{g_buck}'s address for the
    tunables rule. Self-tested: of 24 broken versions of this change, 16
    passed the first eight cases (12 accepting more, 4 rejecting more),
    and cases for them were added. Every broken version, 50 in all, now
    fails the self-tests except the two that cannot behave differently.
    A rule that the
    arming write is window-checked was considered and not added (it needs
    the counter read and both tests located on every path to the write). \\
Results (this build).
  & \texttt{make}: 0 warnings; \texttt{make check} ok; \texttt{make test}
    0 failures, output identical to the build before the fix. Mutants:
    caught by \texttt{make check}: CMP1 after OENR, no OENR, the first
    pass's write unguarded, its CYCCNT read after the \code{strcc}, the
    gains' check, the tunables' copy or the score coefficients after the
    wait; review only (pass all three checks): arming without the window,
    without the ADC1 re-check, with 5120, or with the window's upper bound
    off by one; a refused arm that does not retry (the window's or the
    ADC1 re-check's refusal ending the wait); \code{buck_ctrl_arm} on
    \code{g_buck.tune} instead of \code{c.t}; the entry's \code{asm}
    removed or Ki$\cdot$T after the wait (only the clock count shows it);
    a bad sample while waiting not counted (until 26 October
    \texttt{make check} failed it on false messages, see the
    \texttt{make check} row) or not ending the wait; a STOP
    ignored while waiting; the raw codes not published while waiting, or
    the memory barrier after the wait for the sample removed;
    \code{out.*} not published; the state into \code{g_buck.st};
    $\hat v$ seeded at 0; \code{state} not set; the arming CYCCNT after
    the counter; the stamps stored before the first CMP1 write; the first
    pass's stamps not stored after its CMP1 write, on its trip or on a
    stop; \code{trip_cyc} not stored. The RAM-execution linker script (used by
    no configuration) now overflows by 24\,B (fits with
    \code{-DBUCK_BENCH=0}). \\
Review of the fix (25 October, late).
  & Two further reviews (an emulator of the linked code, and a clock
    counter written independently of the one above) confirmed the build
    and the figures and found wording to correct, corrected in
    \texttt{Inc/buck.h}, \texttt{Inc/buck\_config.h}, comments in
    \texttt{Src/buck.c}, the README and this memo: the 176--220 clocks at
    4096 had been given as the first write's margin (it is 126--144 at
    every arming point); the arming CMP1 write's position, the tighter
    limit on \code{BUCK_ARM_LATEST}, was not stated; the ordinary copies
    had been called the same instructions; \code{trip_cyc} had been called
    the first instruction of every trip branch; a wrapped \code{first_cnt}
    (below 2048, fault~5: the margin row then reads 6401--8448 and the
    cycle formula needs $n > 1$) and $\mathtt{first\_cyc} -
    \mathtt{arm\_cyc}$ with \code{first_*} $=$ \code{0xFFFFFFFF} were not
    explained; the stamps are not stored on a sample timeout (this was in
    the open points only); \code{g_buck.st} is no longer written at
    arming. No code changed: the \texttt{.bin} is byte-identical (sha1
    e5c7798f\ldots), \texttt{make}, \texttt{make check} and
    \texttt{make test} as above. \\
Third review (26 October).
  & A check of \texttt{make check}'s self-tests against broken versions
    of the checker, and of the documents. Fixed in
    \texttt{tools/check\_loop.py}: the false CMP1 messages and the
    self-test gaps (the \texttt{make check} row); every firmware mutant
    caught before is still caught with the same message, and the others
    still pass, the bad sample not counted now included (Results row).
    Corrected or added in the README and this memo: the review-only list;
    the first pass's whole length (item~4); the ordinary passes' nominal
    figures in the README table's convention (147.1 / 150.1 / 154.1 /
    154.1, 7.06 wait clocks); \code{arm()}'s \code{out.g} store between
    CMP1 and OENR; the derived rows before the first start (8448 and 0);
    a third option for the arming write, and the rebuild before the board
    run (below). No firmware file changed: the \texttt{.bin} is
    byte-identical (sha1 e5c7798f\ldots). \\
Left for you (25 October; status 28 October: the fix ran on 27 October,
reported working, no data read; you chose not to add the stamp rows to the
workspace's list; see Section~\ref{sec:followup1028} for the arming write,
the passes and the expected stamp values of the 28 October build).
  & Rebuild Debug in CubeIDE (or relaunch with build-before-launch) and
    flash it first: the Debug ELF that ran on 25 October (built 21:04) is
    the code before the fix, and restarting inside a session started
    with it runs that code again. The new code is loaded when the
    \code{arm_cnt} Live Expressions row evaluates and reads non-zero after
    a start; the workspace's list still had only the 41 old rows (25
    October evening), and the seven new ones are added with the README's
    \texttt{cp} command with CubeIDE closed (not done; on 28 October you
    chose not to add them).
    Run it on the board and read \code{arm_cnt} (about 2100--2650
    expected), $8448 - \mathtt{first\_cnt}$ (about 4000--4600 ticks,
    126--144 clocks; about 3450--4150 with the board's 16--20\,\%),
    $\mathtt{first\_cyc} - \mathtt{arm\_cyc}$ (326--345; about 335--356)
    and \code{busy_max} (150--157 if the model holds). Open points, each
    implemented in the most conservative reading: the first pass as a
    fifth copy with its step chosen at run time and the copy re-chosen
    after it (four copies with the step built in: 1.4\,KB more and 2--3
    clocks later CMP1 writes); a debugger write of \code{state} or
    \code{warmup} during the wait is not seen; the stamps are stored by
    the first pass, so a sample timeout (fault~4) or a CPU fault before
    it leaves the previous start's values, which nothing in \code{g_buck}
    tells apart (storing them at the arming or on the timeout path would
    change the design); \code{0xFFFFFFFF} for a first pass that stopped
    before its check; the ordinary copies' normal passes are not
    instruction-identical to the build before (item~4: whether that is
    acceptable); the stamps' CYCCNT reads sit just before the counter
    reads, on the paths (the first pass's costs 1~clock on its write;
    read right after the \code{strcc} instead, 5--6~clocks after the
    counter read, it would be off the path, but \texttt{make check} then
    needs the \code{bcs} rule widened); \code{BUCK_ARM_LATEST} kept at
    4096 though the arming write lands only 24--37 clocks (pessimistic)
    before 264 there (lower it; store \code{out.*} after OENR: 6--7
    clocks; or hold \code{buck_ctrl_arm}'s float constants across the
    wait, loaded in the entry, not built: between the arming check's
    counter read and that write the board's path makes 10 CCM literal
    loads, 16~clocks static and 27 nominal together, nine of them those
    constants, without whose loads it is 13~clocks static and about 23
    nominal shorter, at a possible cost to the passes' register
    allocation, as the entry's \code{asm} with its values as
    \texttt{"+t"} outputs (counted, not kept) cost 1--3 clocks per pass);
    a window rule in \texttt{make check}; the
    RAM script's 24\,B; two comments in the host test still describe the
    seed and the gains' check at ``the loop entry'' (test unchanged). \\
\bottomrule
\end{xltabular}}

% ============================================================
\section{Follow-up Requests, 28 October (Buck\_disturb\_obs: fewer clocks, the same results)}
\label{sec:followup1028}
% ============================================================

You asked for the loop to use fewer clocks while keeping the trigger
score, the PI calculation, the observer and the sensing as they are,
with bit-identical results only (NaN payloads included; FMA stays off).
Two scouting rounds ran in scratch copies of the project, each candidate
checked against the 25~October evening build (bin e5c7798f\ldots, the
build you reported working on the board on 27~October) in a differential
emulator. On 28~October you chose the ``recommended merged'' build, and it
was applied in place. Its verification found one more operand swap (row~5),
which you had pinned the same day (row~11); the check of the pin corrected the
figure for the read instants of the tunables, and you accepted the measured
instants (row~12). Clock figures are counted on the linked ELF by two
counters, both with your static model and the boot benchmark's wait states
(nominal; pessimistic with CCM literals at $+2$): counter~1 (c1) costs the
paths the differential emulator executes over 26 scenarios, counter~2 (c2)
walks every path. Both are scratch tools, not in the project. Nothing of
this build has run on the board.

{\small
\renewcommand{\arraystretch}{1.25}
\begin{xltabular}{\linewidth}{@{}>{\raggedright\arraybackslash}p{0.27\linewidth}L@{}}
\caption{Requests and decisions, 28 October (the clock optimisation).}
\label{tab:followup1028}\\
\toprule
\textbf{Request or decision} & \textbf{Answer and result} \\
\midrule
\endfirsthead
\toprule
\textbf{Request or decision} & \textbf{Answer and result} \\
\midrule
\endhead
1.\ Fewer clocks; keep the score, the PIs, the observer and the
sensing; bit-identical results only; FMA off.
  & Round~1: path, bus, math and structure scouts, merged into one set
    (SET~I). Round~2: SET~I rebased without P3 and re-searched, the
    post-pass variants (tscout), and their merge. Every candidate ran
    through the differential emulator against the 25~October evening
    build; two defects found on the way (F1, F2 below) were fixed before
    the choice. \\
2.\ Apply the ``recommended merged'' build (your choice, 28~October).
  & Done, after a full backup of the project: the patch (sha1
    3f11d22f\ldots, 1001 lines) changes \code{Src/buck.c},
    \code{Inc/buck_ctrl.h} and \code{Inc/buck_config.h} (a comment); bin
    2f144a53\ldots{} (12\,540\,B; \code{.ccmtext} 9888\,B; \code{run_loop}
    8764\,B), the bin the scouting round checked. Then comments only:
    \code{Src/buck.c}, \code{Inc/buck_ctrl.h}, \code{Inc/buck_config.h},
    \code{Inc/buck.h}, the README and this memo; the bin is unchanged.
    The pin of row~11 followed the same day (bin 5b99a0d5\ldots).
    \code{g_buck} is unchanged (248\,B; \code{busy_max} at offset 156, the
    stamps at 228--244). What changed, each the same operations on the
    same values: the arming writes CMP1 (from $v_\mathrm{out}$ alone) and
    OENR first, everything else after a memory barrier, with the arming's
    and the first pass's constants in registers loaded at the entry; the
    first pass always takes the UP step; the deadline check is an
    \code{asm goto}; the valley and integrator updates follow the CMP1
    write; the pass's I/O is \code{ldr}/\code{str} at offsets from ADC1's
    base; ADC2$\times$0.25 is one \code{vcvt}; BSRR from one MOV and
    \code{out.g} from the BSRR word; a pre-biased Pad\'e table base; the
    gains' check forms both difference orders; the copies' constants in
    registers formed once after the first pass (\code{run_ctx_t.q}); the
    first pass goes straight into the copy of its mode; the telemetry pass
    out of line; \code{busy_max}'s second CYCCNT read through the DWT base
    in a register (BD); the post pass (row~3). \code{RUN_CODE}:
    \code{hot, optimize("O3", "no-if-conversion", "no-crossjumping",
    "no-schedule-insns")}, copies in the order HOLD, UP, DOWN, FREE\_UP
    (before: \code{"no-tree-coalesce-vars", "no-tree-ter"}; UP, FREE\_UP,
    HOLD, DOWN), chosen from 35 flag sets $\times$ 24 orders (840 builds)
    as no worse than the 25~October evening build on any target in either
    counter and the best on the most. \\
3.\ The exact $\beta/\gamma$ check stays in the tunables pass (your
decision; P3, which moved it into the next pass, rejected).
  & Kept there and at the loop entry, unchanged. Only the target of the
    new tunables and the reference step for the next 64 passes moved: the
    pass after a tunables pass (the post pass) runs in the same copy and
    forms them after its CMP1 write, then takes that step itself. The
    values are the same and are used from the same passes. \\
4.\ Rejected options.
  & P3 (the gains' check in the pass after the tunables pass);
    S3 (the first pass as an ordinary pass of its copy: first pass 43
    clocks shorter, but the arming path back at the 25~October length and
    \code{first_cnt}, \code{first_cyc} lost); C10 (the arming CMP1 from a
    16\,KB CCM table); B3/B8 (the tunables pre-written into an appended,
    debugger-writable \code{g_buck} field). \\
5.\ NaN payloads.
  & Fixed by the merge (F1): the current PI's product $k_p e$ in the
    post pass had its operands swapped, which changes the NaN payload in
    \code{g_buck.st.integ_i} at the next stop when both are NaN; its order
    is now pinned (\code{buck_mul_nm}: $k_p$ first in the first pass, $e$
    first in the copies, as the 25~October build), like the five earlier
    pins (current scaling, $\alpha e$, the average, the valley, the score's
    two additions). Your first decision: no new pins; the observer's
    prediction and the clamped PI sum unpinned (three two-NaN sites,
    listed at the end of this row, whose NaN payloads cannot reach an
    output). The check
    after the change
    (row~9) found that
    the prediction $\hat v + a(i_L - \hat\imath)$ had its operands swapped
    in the post pass of every copy ($\hat v$ first; the product first in
    the 25~October build), which changed \code{out.i_load}'s NaN payload
    in one narrow case. Your second decision (row~11): pinned there, the
    product first (\code{buck_obs_step_ao}, with \code{buck_add_nm}); the
    first pass ($\hat v$ first in both builds) and the other passes keep
    the compiler's order. Three two-NaN sites stay unpinned (row~9): the
    voltage-PI and the current-PI clamped sums and the current PI's
    product $k_i T e$, which feeds only its sum. The clamp turns a NaN sum
    into the lower limit, so none of their NaN payloads can reach an
    output. \\
6.\ Tunables read once per 64 passes.
  & Fixed by the merge (F2): GCC re-loaded the 14 tunables from
    \code{g_buck.tune} at the top of the post pass in a
    \code{BUCK_ADC_CSR_POLL = 0} build (legal C, but a debugger write at
    that pass's sample would be taken one pass early). A memory barrier
    right after \code{c->t = b->tune} prevents it; the tunables are read in
    exactly the base's passes (row~9). With
    \code{BUCK_ADC_CSR_POLL = 0} the passes use \code{wait_sample()}
    (P1). \\
7.\ \code{trip_cyc}: comments only (your decision).
  & No code change. Every trip branch, the first pass's included, now
    starts with a literal load of the DWT base from CCM and then the
    CYCCNT read (in this, the \code{BUCK_ADC_CSR_POLL = 1} build; built
    with 0, one branch has a register move between the two): 11~clocks static after the counter read in the ordinary
    copies, 10 in the first pass, whose counter read pipelines behind its
    own CYCCNT read (14 / 13 with the counter read's wait, about 1 more
    nominal for the literal's); before 11--16. Comments in
    \code{Src/buck.c}, \code{Inc/buck.h} and the README updated. \\
8.\ \code{busy_max} and the five stamps kept; \texttt{make check} with
\code{BUCK_ADC_CSR_POLL = 0} left failing; the stamp rows not added to the
workspace's Live Expressions (your decisions).
  & \code{busy_max}'s window is 137 / 137 / 142 / 142 static, 144 / 144 /
    149 / 149 nominal (HOLD / FREE\_UP / UP / DOWN; before 143 / 146 /
    150 / 150 and 150.12--157.12). Its second CYCCNT read no longer
    follows a literal-pool load inside the window (2~clocks static, about
    3 nominal in the 25~October build), so the board reads a few clocks
    less for the same work. The checker fails both
    \code{BUCK_ADC_CSR_POLL = 0} builds (base and new) with ``no pass head
    (ADC12\_CSR read) found''; \texttt{tools/check\_loop.py} is unchanged.
    The workspace's \texttt{.metadata} was not touched. \\
9.\ Proof of bit-identity and its scope.
  & Differential emulator (scratch): both linked ELFs run instruction by
    instruction on the same sample streams and debugger writes; compared:
    every peripheral write and its order within each pass, every
    \code{g_buck} field a debugger can see pass by pass (\code{busy_max},
    \code{bench}, \code{rst_cause} and the stamps reported, not compared),
    faults and their moments. Suites (each with the new build, bin
    5b99a0d5\ldots, at 1.0$\times$, 1.3$\times$ and 0.6$\times$ speed,
    strict write order): 901 scenarios (the standard 887 and 14 around the
    tunables and post passes), 1641 (the 887, 46 at the post pass, 708
    with debugger writes at the post pass's sample, gain sequences and NaN
    pairs per transition), 1069 (the 887 and 182 with NaN tunables taken
    up at a tunables pass and a STOP right after the post pass), and the
    verifier's \texttt{vbit} group, 352 scenarios aimed at the post pass
    (random and special tunables, NaN, $\pm\infty$, subnormal sensing and
    score values, every copy transition, stops and restarts): all the same
    except 5 of the 708 (\code{vt_tr_stall_*}), where the 25~October
    build trips its own fault~5 after a debugger stall and the new one
    does not (the same with the base at 0.9$\times$). Before the pin of
    row~11 (bin 2f144a53\ldots) 14 of the 352 differed at all three
    speeds: when the measured current overflows to $\pm\infty$ in a
    tunables pass (only with extreme sensing tunables written by the
    debugger, e.g.\ a current gain of the order of $10^{35}$ or more with a
    full-scale sample) and the sensing tunables that pass takes are NaN,
    that build published the default NaN in \code{out.i_load} where the
    25~October build publishes the tunable's NaN (the A2 code 0 in both).
    Instruction-level NaN check (\texttt{nanops.py}; every FP operation
    with two NaN operands, in the 611 NaN scenarios of the standard suite
    and in the 352): at every pinned site, the observer's prediction
    included, the base's result in every case there; the only differences
    are at the clamped PI sums (the voltage and the current PI's sum, and
    the current PI's product $k_i T e$, which feeds only its sum: 116 and
    276 cases of pass, operation and operands), which no output can
    carry. That check sees an operand-order
    difference only where two NaNs of different payloads meet. The
    operand-order check (\texttt{opsorder.py}, added by the verification
    after the change) records the operand order of every executed FP
    addition and multiplication whatever the values, in the 901 and the
    352: every pinned site keeps the base's order, the observer's
    prediction now included (before the pin it was swapped in the post
    pass of every copy); the other operations whose order changed cannot
    meet two NaN operands, or their result cannot reach an output.
    Tunables-read check: \code{g_buck.tune} is read in the same passes in
    both builds. \code{BUCK_ADC_CSR_POLL = 0}: 0 warnings, no ADC12\_CSR
    read in \code{run_loop}, the 901 scenarios the same at all three speeds
    against the base built the same way. \texttt{make test} compiles
    \code{Inc/buck_ctrl.h} only; its output is identical, line for line
    (134 lines: the compile command, the test program's name and its 132
    lines of output). Scope: the emulator's scenarios and its peripheral model;
    timing on the board is not modelled beyond the clock counters. Your
    decisions of rows~11 and~12 set the rest of the scope: FPSCR (its cumulative
    exception flags IOC, OFC, \ldots{} and NZCV) is not part of the
    results, since no code reads its flags (the only FPSCR access in the
    code is \code{vmrs APSR_nzcv, fpscr} right after an FP compare, which
    reads the NZCV that compare has just set; there is no
    \code{vmsr}); a debugger write to \code{g_buck.tune} must be taken in
    the same pass as in the base, and it is (the emulator injects writes
    at pass boundaries), while the instants of the reads moved: in the
    tunables pass the 12 fields other than the observer gains 0--47 clocks
    earlier and the gains from 13 earlier to 9 later; at the loop entry
    \code{vref_counts} 35--54 and the gains 22--38 clocks later
    (row~12). \\
10.\ Clock figures.
  & Table~\ref{tab:clocks1028}. Against the 25~October evening build
    the new build is no worse on any target (CMP1 after EOC, normal pass
    worst and mean, \code{busy_max} window, longest pass, first pass and
    its CMP1, arming path, arming write), static, nominal, pessimistic and
    with VLDR/VSTR at 2~clocks (FP2), in either counter. Higher, not
    targets: the post pass (it carries the target and ramp-step work; its
    CMP1 at 68 / 74 after EOC) and the loop entry in the stopped period
    (constants loaded there; the stopped cycle that sees RUN still reaches
    the arming wait at most 200~clocks, pessimistic, after its EOC read).
    Against round~1's SET~I (scratch only) no worse static and nominal,
    but not under FP2: every normal pass $+1$ (151 against 150, c1), and
    FREE\_UP's \code{busy_max} window $+1$ static (144 against 143 by c1,
    143 against 142 by c2; nominal 151 against 151.06, c1); HOLD's window
    is 1 shorter and UP's and DOWN's are equal, so the worst copy's window
    is equal. The pin of row~11 changes none of these figures. \\
11.\ Your decisions of 28~October, second set (after the verification of
the applied build).
  & (1) Pin the observer's prediction with the product first in the post
    pass only. The verifier found (BITID-1, reproduced twice) that the post
    pass of every copy added it with $\hat v$ first, the only place where
    the compiler chose that order in the copies. Done:
    \code{buck_obs_step_ao} in \texttt{Inc/buck\_ctrl.h}, called with
    \code{epoch != 0} in \code{run_pass()}; the ordinary passes and the
    first pass keep their instructions. Bin 5b99a0d5\ldots{} (the same
    sizes); only the post pass's instructions changed (the four
    \code{vadd}s swap their operands, and the 4 to 11 instructions after
    each are reordered: most move by one place; in the UP, DOWN and
    FREE\_UP copies one multiplication moves two places and one
    subtraction five, three in FREE\_UP; no store or peripheral access
    moves); both counters' outputs are byte-identical
    to 2f144a53's, FP2 included (0~clocks); results in row~9. Your
    first-set decision ``no new pins'' rested on a wrong premise, which I
    gave you with the options: that the instruction-level NaN check guards
    the unpinned sites. That check sees a swap only where two NaNs of
    different payloads meet in one of its scenarios; the operand-order
    check records the operand order of every FP addition and
    multiplication the scenarios execute, NaN or not, and compares it with
    the base's pass by pass, so it sees a swap at every such instruction
    that both builds execute in the same pass with two different operands,
    and it showed this one. It never compares eight post-pass
    instructions, which run one pass later than in the base (it formed them in the tunables
    pass): the ramp target's multiplication by \code{BUCK_V_PER_COUNT}
    and \code{buck_ramp_mode}'s addition of \code{BUCK_RAMP_FAR_V}, in
    each of the four copies. Each has a constant operand, so a swap there
    cannot change a NaN payload, and the conclusion holds.
    (2) FPSCR, its cumulative exception flags (IOC,
    OFC, \ldots) and NZCV, is outside ``results'': no code reads its flags
    (row~9). The new build's gains' check forms both difference orders
    unconditionally, so it can set IOC or OFC where the base sets none.
    (3) Debugger writes to \code{g_buck.tune}: taken in the same pass as
    before is the requirement; the instants of the reads inside the
    tunables pass may move. The figure I gave you with this question,
    ``2--37 clocks earlier'', was wrong, and I did not mention that at the
    loop entry \code{vref_counts} and the gains are read later; your third
    set (row~12)
    accepts the measured instants, the loop entry's included (in the
    tunables pass from 47 clocks earlier to 9 later, at the loop entry
    from 28 earlier to 54 later); a write that lands between the two
    builds' reads of a field is taken one reload apart. (4) The
    RAM-execution linker script's overflow (P8: 3096\,B; it no longer fits
    with \code{-DBUCK_BENCH=0}) is left as it is. \\
12.\ Your decisions of 28~October, third set (after the check of the
pin).
  & (1) The read instants of \code{g_buck.tune}: you accepted the
    measured instants, the loop entry's included; documents only, no code
    change (not taken: reading the tunables earlier at the loop entry, a
    code change, not explored). The figure given to you before, ``2--37
    clocks earlier'' for the 12 fields other than the observer gains in
    the tunables pass (row~11 (3), and in row~9 and the README until
    then), was wrong: it was measured on one scenario and does not
    reproduce there (\code{plain}: in its ten tunables passes those fields
    5--42 clocks earlier, 5--31 in pass~87). Measured over the 901 scenarios (the first
    load of each field, in clocks after the pass's EOC, at 1.0$\times$;
    the 25~October evening build against bin 5b99a0d5\ldots): in the
    tunables pass the 12 fields other than the observer gains are read
    0--47 clocks earlier (each field's 1st to 99th percentile within
    3--45 earlier) and the
    gains from 13 earlier to 9 later. At the loop entry (the stopped
    cycle that sees RUN, which a debugger halt can delay, as in the
    scenario \code{halt_in_entry})
    \code{vref_counts} is read 35--54 clocks later, the gains 22--38
    later, \code{kp_v} from 5 earlier to 14 later and the other ten from
    28 earlier to 9 later. \code{g_buck.tune} is read in the same passes
    in both builds (13\,285 passes, no scenario differs). A write that
    lands between the two builds' reads of a field is taken one reload
    apart, by the build that reads that field first only at its next
    reload: in the tunables pass that is mostly the new build (it takes
    the write a reload later), at the loop entry, for \code{vref_counts}
    and the gains, the 25~October build (the new build takes it a reload
    earlier). A write that spans several fields can be split between old
    and new values differently. (2) The comment-only edits beyond the pin
    are kept: in \code{Src/buck.c} ``(27 Oct)'' became ``(scouted 27 Oct,
    applied 28 Oct)'' at five places; in \code{Inc/buck_ctrl.h} the
    comments on \code{buck_iload_obs()} no longer say that
    \code{run_pass()} calls it (the host tests do; \code{run_pass()} runs
    the same step through \code{buck_obs_step_ao} and
    \code{buck_iload_dac_k}), and the pin's comment says ``so did the
    28 Oct build before this pin (bin 2f144a53)''. They change no
    instruction (bin 5b99a0d5\ldots). Corrected at the same time,
    documents only: row~11 (1) (the eight instructions the operand-order
    check never compares), the line count of \texttt{make test} (row~9)
    and the unpinned clamped PI sums (row~5 and below). \\
Left open.
  & Not run on the board: rebuild Debug in CubeIDE before flashing
    (CubeIDE rebuilt \texttt{Debug/} at 21:33 on 28~October; its ELF
    gives bin 5b99a0d5\ldots; no board run with it has been reported).
    Expected on the board (estimates from the counters, no board data):
    \code{arm_cnt} about 1800--2350 ticks; $8448 - \mathtt{first\_cnt}$
    about 4700--5250 ticks (146--164 clocks; about 4250--4900 with the
    board's 16--20\,\%); $\mathtt{first\_cyc} - \mathtt{arm\_cyc}$ about
    315--334 (322--343); \code{busy_max} about 144--149 (167--179).
    Left as it is (your decision, row~11): P8, the RAM-execution linker
    script's overflow, now 3096\,B (101\,400 of 98\,304\,B; 2072\,B with
    \code{-DBUCK_BENCH=0}, which no longer fits; 24\,B before).
    Kept as they are, unanswered: the wiring checks (V3-1: wiring mistakes
    in \code{Src/buck.c} that pass all three checks, review only). The
    clamped PI sums stay unpinned (row~5: three two-NaN sites, the
    voltage-PI and the current-PI sums and the current PI's product
    $k_i T e$, which feeds only its sum; none can carry its NaN payload
    to an output); \texttt{make check} fails with
    \code{BUCK_ADC_CSR_POLL = 0} (row~8). The scratch timing tool
    \texttt{vtime.py} has a wrong comment on \code{g_buck.out}'s offset
    (\code{out} is at 0x44, \code{st} at 0x70); it changes no figure. \\
\bottomrule
\end{xltabular}}

\begin{table}[H]
\centering
\caption{Clock figures, 28 October build (bin 5b99a0d5\ldots; 2f144a53\ldots{}
the same) against the
25 October evening build (in brackets), static / nominal. ``Table'' is
the README convention (body plus the closing status read); ``with head''
from the status read that finds EOC to the next pass's first status read.
c1 = c2 where one value is given.}
\label{tab:clocks1028}
\footnotesize
\renewcommand{\arraystretch}{1.15}
\begin{tabularx}{\linewidth}{@{}>{\raggedright\arraybackslash}p{0.34\linewidth}LL@{}}
\toprule
\textbf{Figure} & \textbf{28 October} & \textbf{[25 October evening]} \\
\midrule
CMP1 after EOC, every pass but the first, worst
  & 68 / 74 & c1 82 / 88; c2 79 / 85 (poll head), 82 / 88 (late head) \\
normal pass (table), HOLD / FREE\_UP / UP / DOWN
  & 137 / 138 / 143 / 143; 145 / 146 / 151 / 151
  & 140 / 143 / 147 / 147; 147.06 / 150.06 / 154.06 / 154.06 \\
normal pass, mean of the copies
  & 140.25 / 148.25 & 144.25 / 151.31 \\
\code{busy_max} window
  & 137 / 137 / 142 / 142; 144 / 144 / 149 / 149
  & 143 / 146 / 150 / 150; 150.12 / 153.12 / 157.12 / 157.12 \\
telemetry pass (table)
  & 171 / 172 / 178 / 178; 179 / 180 / 186 / 186
  & 182 / 185 / 189 / 189; 190.12 / 193.12 / 197.12 / 197.12 \\
longest pass with head (UP/DOWN tunables pass); pessimistic
  & 210 / 220.12; 222 & 271 / 294; 309 \\
pass after a tunables pass, with head (not a target)
  & c1 175 / 187.25; c2 177 / 189.25 & an ordinary pass, 154 / 161.06 \\
first pass, EOC to the next head; its CMP1 after EOC
  & 188 / 204.5; 68 / 74
  & c1 219 / 248.31, c2 221 / 250.31; c1 82 / 95.44, c2 81 / 94.44 \\
arming counter read to the first pass's guarded counter read
  & c1 144 / 153.12; c2 141 / 150.12 & c1 194 / 220.19; c2 195 / 221.19 \\
arming counter read to the arming CMP1 write
  & 29 / 33 & 88 / 102.62 \\
loop entry, stopped period (not a target)
  & c1 120 / 139.12; c2 124 / 145.25 & c1 103 / 111.5; c2 103 / 113.62 \\
worst normal pass with VLDR/VSTR at 2 clocks (FP2)
  & c1 151 / 159; c2 150 / 158 & c1 156 / 163.06; c2 154 / 161.06 \\
\bottomrule
\end{tabularx}
\end{table}

% ============================================================
\section{Follow-up Request: Live Expressions}
% ============================================================

\begin{description}
\item[Request.] Clear the old Live Expressions and provide a new list
for this firmware.
\item[Result.] The old 44-row list (names from
\code{Buck_Voltage_Mode}) was removed from the CubeIDE workspace and
backed up as \code{saved_expr.dat.pre_current_mode_bak}. A new
35-row list was installed, and a copy is kept in
\code{tools/live_expressions.txt}. On 27 September
(Section~\ref{sec:followup}) \code{bench}, \code{tune.ev_norm},
\code{tune.du_norm} and \code{trip_cnt} were added, and
\code{tune.iload_n} was added and later removed: 39 rows. On 28 September
(Section~\ref{sec:followup28}) \code{tune.iload_s0} and
\code{out.i_load_lvl} were added: 41 rows in all. In BUCK\_Adpt\_Learning
(24 October) they are replaced by \code{tune.iload_tau0} and
\code{out.i_load_alpha}: still 41 rows. On 25 October \code{tune.rho},
\code{tune.ev_norm} and \code{tune.du_norm} were replaced by
\code{tune.score_a}, \code{tune.score_b}, \code{tune.score_c} and
\code{tune.iload_s_th}: 42 rows. In this copy, Buck\_disturb\_obs
(25 October), the two observer gains \code{tune.obs_beta} and
\code{tune.obs_gamma} replace the three rows \code{tune.iload_tau0},
\code{tune.iload_s_th} and \code{out.i_load_alpha}: 41 rows. On the
evening of 25 October (Section~\ref{sec:followup1025f}) the start's
timing stamps \code{arm_cnt}, \code{arm_cyc}, \code{first_cnt},
\code{first_cyc}, \code{trip_cyc} and the two rows
\code{first_cyc-arm_cyc} and \code{8448-(int)first_cnt} were added:
48 rows. You chose (28 October) not to add these seven rows to the
workspace's Live Expressions list; they are in
\code{tools/live_expressions.txt}, and the command below restores the
whole list if wanted. Every row was checked with
\texttt{gdb} against both the command-line and the CubeIDE Debug build
(Buck\_disturb\_obs: against the command-line build).
\end{description}

\begin{table}[H]
\centering
\caption{Groups in the new Live Expressions list.}
\label{tab:live}
\small
\renewcommand{\arraystretch}{1.25}
\begin{tabularx}{\linewidth}{@{}l r L@{}}
\toprule
\textbf{Group} & \textbf{Rows} & \textbf{Contents} \\
\midrule
Controls    & 18 & \code{command}, \code{state}, \code{fault};
                     all \code{tune.*} values (V$_{\mathrm{ref}}$ in counts
                     and volts, four PI gains, $\alpha_I$,
                     $S_{\mathrm{th}}$, the score coefficients $a$, $b$,
                     $c$, the observer gains $\beta$ and $\gamma$, A1 zero
                     and gain) \\
Raw ADC     & 3  & $v_{\mathrm{out}}$ in V, A1 in counts, A1 current in A;
                     live even when stopped, for the A1 zero calibration \\
Outputs     & 13 & \code{out.*}: $v_{\mathrm{out}}$, $V_{\mathrm{ref}}$,
                     $e_v$, measured and estimated average current,
                     residual, peak, the observer's load current,
                     $u[k]$, $S[k]$, $g[k]$,
                     CMP1 and the physical duty; live only while running \\
Diagnostics & 14 & \code{busy_max} (keep below about 240),
                     \code{loops}, \code{adc_miss}, \code{warmup},
                     reset cause, bus benchmark \code{bench},
                     \code{trip_cnt} and \code{trip_cyc} (fault 5), the
                     start's stamps \code{arm_cnt}, \code{arm_cyc},
                     \code{first_cnt}, \code{first_cyc}, and
                     \code{first_cyc-arm_cyc},
                     \code{8448-(int)first_cnt} \\
\bottomrule
\end{tabularx}
\end{table}

\noindent
CubeIDE overwrites its copy when the view or the IDE closes. To restore
the list, close CubeIDE and run from the project folder:
\begin{center}
{\footnotesize\code{cp tools/live_expressions.txt ../.metadata/.plugins/com.st.stm32cube.ide.mcu.livewatch/saved_expr.dat}}
\end{center}

\end{document}
